
\clearpage
\numberwithin{equation}{section}
\setcounter{table}{0}
\renewcommand{\thetable}{O.\arabic{table}}
\setcounter{figure}{0}
\renewcommand{\thefigure}{O.\arabic{figure}}
\setcounter{footnote}{0} 
\setcounter{proposition}{0}
\renewcommand{\theproposition}{O.\arabic{proposition}}
\appendix
\onehalfspacing
% \singlespacing
% \newgeometry{left=1in,bottom=1in, top=1in, right=1in}
\part*{Supplementary Appendix}\label{OA_appendix}

\setcounter{page}{1}
\renewcommand{\thesection}{O\arabic{section}}
\section{The case with an arbitrary number of states}\label{app:Nstates}

\subsection{Environment}

We consider an extension of the model in Section \ref{sec:environment.sec} where aggregate productivity growth $x_{t}$ takes $N$ possible values, i.e. $x_t \in \{x^1, x^2, \ldots, x^N\} \equiv \mathcal{X}$, where $x^1 < x^2 < \ldots < x^N$.
The objective probability of switching from state $s \in\{1,2,\ldots, N\}\equiv \mathcal{S}$ to state $s'\in \mathcal{S}$ is denoted by $p_{ss'}$ and the corresponding subjective probability for household $i\in \mathcal{I}$ is denoted by $p_{ss'}^i$. 
Households can trade Arrow securities that pay off conditional on every possible state. 
We also assume that households die with probability $\kappa$ and leave their financial wealth to their child, which will have type $j$ with probability $\mu_j$. 
This assumption ensures that a non-degenerate stationary distribution of wealth exists.
The next proposition provides a characterization of the equilibrium in this $N$-state economy. The main conclusion is that the results of Section \ref{sec:environment.sec} are essentially unchanged in this more general setting.

\begin{proposition}[N-state economy]\label{prop: oa_Nstate} Suppose that $x_t \in \mathcal{X}$, where $x_t$ takes $N$ possible values.

\begin{enumerate}
    \item The (scaled) household's problem can be written as follows
    \begin{equation}
    \frac{v_i(X,s)^{1-\psi^{-1}}-1}{1-\psi^{-1}} = \max_{c_{i},R_{i,n}'}  (1-\beta) \frac{c_i^{1-\psi^{-1}}-1}{1-\psi^{-1}} + \beta \frac{\mathbb{E}_i\left[ (v_i(X',s')n')^{1-\gamma}\right]^{\frac{1-\psi^{-1}}{1-\gamma}}-1}{1-\psi^{-1}},
\end{equation}
subject to the flow budget constraint $n' = R_{i,n}' (1 - c_i) $, the natural borrowing limit $n' \geq 0$, and the portfolio-return constraint
\begin{equation}
    \mathbb{E}_s[\Lambda' R_{i,n}'] = 1.
\end{equation}

\item Consumption-wealth ratio and the investor's SDF are given by 
\begin{align}
    c_i(X,s) &= \frac{ (\beta^{-1}-1)^{\psi} \mathcal{R}_i(X,s)^{1-\psi}}{1+ (\beta^{-1}-1)^{\psi} \mathcal{R}_i(X,s)^{1-\psi}},\\
    \Lambda_i(X,s,s') &=   \beta^\theta \left(\frac{c_i(\chi(X,s,s'),s') N'}{c_i(X,s) N}\right)^{-\frac{\theta}{\psi}} R_{i,n}(X,s,s')^{-(1-\theta)},
\end{align}
and the change-of-measure condition is given by
\begin{equation}
    \Lambda_i(X,s,s') = \frac{p_{ss'}}{p_{ss'}^i} \Lambda(X,s,s').
\end{equation}

\item Wages, hours, and profits are given by
\begin{equation}
    h(E) = \left( \frac{\alpha E}{\xi} \right)^{\frac{1}{1+\nu -\alpha}}, \qquad  
    w(E) = \xi \left( \frac{\alpha E}{\xi} \right)^{\frac{\nu}{1+\nu -\alpha}}, \qquad  
    \pi(E,s) = \left( \frac{\alpha}{\xi} \right)^{\frac{\alpha}{1+\nu-\alpha}} \left[ x_{s} E^{\frac{\alpha}{1+\nu-\alpha}}  - \alpha E^{\frac{1+\nu}{1+\nu-\alpha}}  \right].
\end{equation}
\item The law of motion of the endogenous aggregate state variables is given by
\begin{align}
    \mathcal{L}'(X,s) &= \sum_{s'\in \mathcal{S}} \frac{p_{ss'} \Lambda(X,s,s')}{\sum_{\tilde{s} \in \mathcal{S}} p_{s\tilde{s}} \Lambda(X,s,\tilde{s})} x_{s'}, \\
        \eta_i'(X,s,s') &= \frac{(1-\kappa) \eta_i R_{i,n}(X,s,s') (1-c_i(X,s))}{\sum_{j=1}^I\eta_j R_{j,n}(X,s,s') (1-c_j(X,s)) } + \kappa \mu_i.
\end{align}

\item The market clearing conditions for consumption and the Arrow security for state $s\in \mathcal{S}$ are given by
\begin{equation}
\sum_{i=1}^I \eta_i c_i(X,s) = \frac{ x_s h(E)^\alpha - \xi \frac{h(E)^{1+\nu}}{1+\nu}}{P(X,s)}, \qquad \qquad   \sum_{i=1}^I \tilde{\eta}_i R_{n,i}(X,s,s') = R_p(X,s,s'),
\end{equation}
where $\tilde{\eta}_i \equiv \frac{\eta_i (1-c_i(X,s)}{\sum_{j=1}^I \eta_j (1-c_j(X,s)}$.
\end{enumerate}
\end{proposition}

\begin{proof}
See Online Appendix \ref{proof: oa_Nstate}.
\end{proof}

An implication of the result above is that the LDF corresponds to the risk-neutral expectation of productivity growth. 
The following corollary shows that $\mathcal{L}'(X,s)$ can be expressed as the expected productivity growth (under the objective probability measure) discounted by a risk premium.

\begin{corollary} Let $R_g(X,s,s')$ denote the return on a claim on productivity growth, then
\begin{equation}
    \log \mathcal{L}'(X,s) = \log \mathbb{E}_s[x_{s'}] - \log \overline{R}^e_{g}(X,s),
\end{equation}
where $\overline{R}^e_{g}(X,s) \equiv \sum_{s'\in \mathcal{S}} p_{ss'} \frac{R_{g}(X,s,s')}{ R_b(X,s)} $ is the risk premium on a claim on productivity growth.
\begin{proof} The price of a claim on productivity growth is given by
\begin{equation}
    P_g(X,s) = \mathbb{E}_s[ \Lambda(X,s,s') x_{s'} ],
\end{equation}
and the return on this claim is given by $R_g(X,s,s') = \frac{x_{s'}}{P_g(X,s)}$.

Expressing the pricing condition above in terms of covariances, we obtain
\begin{equation}
    \mathbb{E}_s[R_g(X,s,s')] - R_b(X,s) = - Cov_s\left( \frac{\Lambda(X,s,s')}{\mathbb{E}_s[\Lambda(X,s,s')]} ,  \frac{x_{s'}}{P_g(X,s)} \right).
\end{equation}

Similarly, we can write $\mathcal{L}'(X,s)$ in terms of covariances:
\begin{equation}
    \mathcal{L}'(X,s) = \mathbb{E}_s[x_{s'}] + Cov_s\left( \frac{\Lambda(X,s,s')}{\mathbb{E}_s[\Lambda(X,s,s')]} ,  x_{s'}\right).
\end{equation}

Using the fact that $P_g(X,s) = \mathcal{L}'(X,s) / R_b(X,s)$, we can combine the expressions above to obtain
\begin{equation}
    \mathcal{L}'(X,s) = \mathbb{E}_s[x_{s'}] -\left(\frac{\mathbb{E}_s[R_g(X,s,s')]}{R_b(X,s)} - 1\right)\mathcal{L}'(X,s) \Rightarrow \mathcal{L}'(X,s) = \frac{\mathbb{E}_s[x_{s'}]}{\overline{R}_g^e(X,s)}.
\end{equation}
    
\end{proof}
    
\end{corollary}

\subsection{Special Case I: Log utility}

We consider next the special case where $\psi = \gamma = 1$ for the economy with an arbitrary number of states. 
Proposition \ref{oa_log_utility} below shows that the main implications from Section \ref{sec:log_utility.sec} extends to this more general economy. 

\begin{proposition}[Log-utility]\label{oa_log_utility} Suppose $\psi = \gamma = 1$ and that the following condition is satisfied $x^N < \frac{x^1}{\alpha}$.

\begin{enumerate}
    \item Consumption and portfolio decisions are given by
    \begin{equation}
        c_i(X,s) = 1-\beta, \qquad \qquad R_{i,n}(X,s,s') = \frac{p_{ss'}^i}{p_{ss'}(X)} R_p(X,s,s').
    \end{equation}
    \item The economy's SDF is given by
    \begin{equation}
        \Lambda(X,s,s') = \frac{p_{ss'(X)}}{p_{ss'}} R_p'(X,s,s')^{-1}.
    \end{equation}
    \item The price and return on the surplus claim are given by
    \begin{align}
        P(X,s) &= \frac{x_s h(E)^\alpha - \xi \frac{h(E)^{1+\nu}}{1+\nu}}{1-\beta}, \\
        R_p(X,s,s') &= \frac{x_s}{\beta} \frac{x_{s'} \mathcal{L}'(X,s)^\frac{\alpha}{1+\nu-\alpha} -\frac{ \alpha}{1+\nu} \mathcal{L}'(X,s)^{\frac{1+\nu}{1+\nu-\alpha}}}{x_{s} E^\frac{\alpha}{1+\nu-\alpha} -\frac{ \alpha}{1+\nu} E^{\frac{1+\nu}{1+\nu-\alpha}}}.
    \end{align}
    \item The risk premium on the surplus claim and the interest rate are given by
    \begin{align}
        R_b(X,s) &= \left(1 - \frac{\alpha}{1+\nu} \right) \frac{x_s}{\beta} \frac{  \mathcal{L}'(X,s)^{\frac{1+\nu}{1+\nu-\alpha}}}{x_{s} E^\frac{\alpha}{1+\nu-\alpha} -\frac{ \alpha}{1+\nu} E^{\frac{1+\nu}{1+\nu-\alpha}}},\\
        \mathbb{E}_s[R_p^e(X,s,s')] &= \frac{x_s}{\beta} \frac{[\mathbb{E}_s[x_{s'}]-\mathcal{L}'(X,s)] \mathcal{L}'(X,s)^\frac{\alpha}{1+\nu-\alpha} }{x_{s} E^\frac{\alpha}{1+\nu-\alpha} -\frac{ \alpha}{1+\nu} E^{\frac{1+\nu}{1+\nu-\alpha}}}.
    \end{align}
    
\item The law of motion of the aggregate state variables are given by
\begin{align}
    \mathcal{L}'(X,s) &= \sum_{s'\in \mathcal{S}} x_{s'} \frac{p_{ss'}(X) \left[x_{s'} -\frac{ \alpha}{1+\nu} \mathcal{L}'(X,s) \right]^{-1}}{\sum_{\tilde{s}\in \mathcal{S}} p_{s\tilde{s}}(X)\left[x_{\tilde{s}} -\frac{ \alpha}{1+\nu} \mathcal{L}'(X,s) \right]^{-1}}\\
    \eta_i'(X,s,s') &= (1-\kappa) \eta_i \frac{p_{ss'}}{p_{ss'}(X)} + \kappa \mu_i,
\end{align}
and there exists a unique value of $\mathcal{L}'(X,s) \in (x_1,x_N)$ satisfying the law of motion of $\mathcal{L}$.
\end{enumerate}
\end{proposition}

\begin{proof}
See Online Appendix \ref{proof:oa_log_utility}.
\end{proof}

\subsection{Special Case II: Representative Agent with IID Returns}

We consider next a different special case which is also particularly tractable: investors have common iid beliefs, $p_{ss'}^i = p_{s'}^*$, and the supply and demand of labor converge to zero. Formally, we assume $\alpha = \hat{\alpha} \epsilon$ and $\xi = \hat{\xi} \epsilon$ and take the limit as $\epsilon$ goes to zero. 
For simplicity, we focus on the case $\kappa = 0$. 
Because labor is chosen in advance, returns on financial assets would not be iid even if the process for aggregate productivity is iid. 
By taking the limit as supply and demand goes to zero, we ensure that all equilibrium objects are well-defined in the limit and the economy behaves essentially as an endowment economy, analogous to an iid version of the \citet{MP85} economy.

Proposition \ref{oa_iid} provides a characterization of this limit economy. To highlight these results apply to this particular limit, we denote the equilibrium objects in the limiting economy with an $*$, e.g. $v^*(X,s)$ and $c^*(X,s)$.


\begin{proposition}[IID Returns]\label{oa_iid} Suppose $p_{ss'}^i = p_{s'}^*$, $\alpha = \hat{\alpha} \epsilon$, and $ \xi = \hat{\xi} \epsilon$. Suppose also the following condition  is satisfied: $\beta^* \equiv \beta \mathbb{E}^*[x_{s'}^{1-\gamma}]^{\frac{1-\psi^{-1}}{1-\gamma}} < 1$. Then, the economy in the limit as $\epsilon \rightarrow 0$ satisfies the following conditions:
\begin{enumerate}
    \item Consumption and portfolio decisions:
    \begin{equation}
        c_i^*(X,s) = 1-\beta^*, \qquad \qquad R_{i,n}^*(X,s,s') = R^*_p(X,s,s').
    \end{equation}
    
    
    \item The net-worth multiplier $v_i^*(X,s)$ is given 
    \begin{equation}
        v_i^*(X,s) = (1-\beta)^{\frac{1}{1-\psi^{-1}}} (1-\beta^*)^{-\frac{\psi^{-1}}{1-\psi^{-1}}}.
    \end{equation}
    
    \item Wages, hours, and profits are given by
\begin{equation}
    h^*(E) = \left( \frac{\hat{\alpha} E}{\hat{\xi}} \right)^{\frac{1}{1+\nu}}, \qquad  
    w^*(E) = 0, \qquad  
    \pi^*(E,s) = x_s.
\end{equation}
    
    \item The economy's SDF is given by
    \begin{equation}
        \Lambda^*(X,s,s') = \beta \mathbb{E}^*[x_{s'}^{1-\gamma}]^{\frac{\gamma-\psi^{-1}}{1-\gamma}} x_{s'}^{-\gamma}.
    \end{equation}
    
    \item The price and return on the surplus claim are given by
    \begin{equation}
        P^*(X,s) = \frac{x_s}{1-\beta^*}, \qquad \qquad 
        R^*_p(X,s,s') = \frac{x_s'}{\beta^*}.
    \end{equation}
    
    \item The risk-free rate and the expected return on the surplus claim are given by
    \begin{align}
        R_b^*(X,s) &= \frac{1}{\beta^*} \frac{\mathbb{E}^*[ x_{s'}^{1-\gamma} ]}{\mathbb{E}^*[x_{s'}^{-\gamma}]} ,  \\
        \mathbb{E}^*[R_p(X,s,s')] &= 
         R_b(X,s) \frac{\mathbb{E}^*[x_{s'}] \mathbb{E}^*[x_{s'}^{-\gamma}]}{\mathbb{E}^*[x_{s'}^{1-\gamma}] }.
    \end{align}
    
    \item The law of motion of the state variables are given by
    \begin{equation}
        \mathcal{L}'(X,s) = E^*, \qquad \qquad \eta_i'(X,s,s') = \eta_i.
    \end{equation}
    where $E^* \equiv \frac{\mathbb{E}^*[x_{s'}^{1-\gamma}]}{\mathbb{E}^*[x_{s'}^{-\gamma}]}$.
\end{enumerate}
\end{proposition}

\begin{proof}
See Online Appendix \ref{proof: oa_iid}.
\end{proof}



The following corollary shows that we recover the standard asset pricing formulae for iid economies in continuous time if we assume that $x_s$ is approximately log-normal.

\begin{corollary} Suppose $\log x_s$ can be approximated by a normal distribution with mean $\mu$ and variance $\sigma^2$. Then, under the assumption of Proposition \ref{oa_iid}, we obtain
\begin{enumerate}
    \item Interest rate:
\begin{equation}
    \log R_b^*(X,s) \approx \rho + \psi^{-1} \left( \mu + \frac{\sigma^2}{2} \right) - \frac{\gamma(1+\psi^{-1})}{2}  \sigma^2,
\end{equation}
where $\rho \equiv - \log \beta$. 

\item Risk-premium:
\begin{equation}
\log \mathbb{E}\left[\frac{R^*_p(X,s,s')}{R_b^*(X,s)} \right]  \approx \gamma \sigma^2.
\end{equation}

\item Risk-neutral expectation of productivity growth:
\begin{equation}
    \log \frac{\mathcal{L}'(X,s)}{\mathbb{E}[x_{s'}]} \approx -\gamma \sigma^2.
\end{equation}
\end{enumerate}

\end{corollary}

The corollary above shows how $\mathcal{L}'(X,s)$ depends on $x_{s'}$ and the equity risk premium. 

\section{Approximate Solution of the General Economy}

In the previous section, we derived exact analytical solutions for two special cases: i) log-utility; ii) homogeneous beliefs and iid returns. 
In this section, we derive asymptotic closed-form solutions for a general economy with an arbitrary number of states, an arbitrary number of households with heterogeneous beliefs, and Epstein-Zin preferences with unrestricted EIS and risk aversion.
The derivations for the benchmark case with homogeneous beliefs and iid returns will be useful in deriving the approximate solution.

\subsection{Perturbation}

Consider a family of economies indexed by $\epsilon$. 
The parameter $\epsilon$ controls three dimensions through which these economies differ from each other. First, it determines the degree of belief heterogeneity:
\begin{equation}
    p_{ss'}^i = p_{s'}^* + \delta_{ss'}^i \epsilon,
\end{equation}
where $\sum_{s' \in \mathcal{S}}\delta_{ss'}^i = 0$. 
We also assume that the objective measure coincides with beliefs in the reference economy, i.e. $p_{ss'} = p_{s'}^*$. Second, $\epsilon$ scales both supply and demand for labor:
\begin{equation}
    \xi = \hat{\xi} \epsilon, \qquad \qquad  \alpha = \hat{\alpha} \epsilon.
\end{equation}
The economy satisfying $\epsilon = 0$ is essentially an endowment economy with iid common beliefs, a special case of the \citet{MP85} economy, as described above.
Third, we assume that $\kappa = \hat{\kappa} \epsilon$, such that there is no mortality risk in the benchmark economy.

All equilibrium objects are now indexed by $\epsilon$. For instance, the net worth multiplier is now given by $v_i(X,s; \epsilon)$. We are interested in the expansion of $v_i(X,s; \epsilon)$ on $\epsilon$, for $\epsilon$ small:
\begin{equation}
    v_i(X,s;\epsilon) = v_i^*(X,s) + \hat{v}_i(X,s) \epsilon + \mathcal{O}(\epsilon^2),
\end{equation}
where $v_i^*(X,s) \equiv v_i(X,s; 0)$ and $\hat{v}_i(X,s)$ represents the first-order correction of $v_i(X,s;\epsilon)$ in $\epsilon$.

Similarly, we can write the consumption-wealth ratio  $c_i(X,s;\epsilon)$ as follows:
\begin{equation}
    c_i(X,s;\epsilon) = c_i^*(X,s) + \hat{c}_i(X,s) \epsilon + \mathcal{O}(\epsilon^2),
\end{equation}
and analogously for the remaining equilibrium variables. 

The functions $v^*_i(X,s)$ and $c^*_i(X,s)$ are already known, as they correspond to the solution of the case with homogeneous beliefs and iid returns, which we characterized above. It remains to solve for $\hat{v}_i(X,s)$,  $\hat{c}_i(X,s)$, and the first-order correction for the other variables. 

We start by providing a characterization of the households' problem in this general economy. 
\begin{proposition}\label{oa_investor_perturbation} Suppose that $p_{ss'}^i = p_{s'}^* + \delta_{ss'}^i \epsilon$, $\alpha = \hat{\alpha} \epsilon$, and $\xi = \hat{\xi} \epsilon$. Suppose also that $\beta^* < 1$. Then,
\begin{enumerate}
    \item Net-worth multiplier:
\begin{equation}
    \frac{\hat{v}_i(X,s)}{v_i^*(X,s)}  =  \beta^*  \sum_{s'\in \mathcal{S}} \omega^*_{s'}\left[\frac{1}{1-\gamma} \frac{\delta_{ss'}^i}{p_{s'}^*} -\frac{\hat{\Lambda}(X,s,s')}{\Lambda^*(X,s,s')} \right] +\beta^* \overline{v},
\end{equation}
where $\omega_{s'}^* \equiv\frac{p_{s'}^* x_{s'}^{1-\gamma}}{\sum_{\tilde{s} \in \mathcal{S}} p_{\tilde{s}}^* x_{\tilde{s}}^{1-\gamma}}$,$X^* = (E^*, \{\eta_i\}_{i=1}^I)$, and
\begin{equation}
    \overline{v} \equiv \frac{ \beta^* }{1-\beta^*}\sum_{\tilde{s}\in \mathcal{S}}  \omega^*_{\tilde{s}}  \sum_{\tilde{s}'\in \mathcal{S}} \omega^*_{\tilde{s}'}\left[\frac{ 1}{1-\gamma} \frac{\delta_{\tilde{s}\tilde{s}'}^i}{p_{\tilde{s}'}^*} -\frac{\hat{\Lambda}(X^*,\tilde{s},\tilde{s}')}{\Lambda^*(X^*,\tilde{s},\tilde{s}')}  \right].
\end{equation}
    
    \item Consumption-wealth ratio:
    \begin{equation}
    \frac{\hat{c}_i(X,s)}{c^*(X,s)} =  (1-\psi) \frac{\hat{v}_i(X,s)}{\hat{v}^*(X,s)}.
\end{equation}

    \item Portfolio return:
\begin{equation}
    \frac{\hat{R}_{n,i}(X,s,s')}{R^*_p(X,s,s')} = \frac{1}{\gamma}  \text{myopic}  + \frac{1-\gamma}{\gamma} \text{hedging},
\end{equation}
where 
\begin{align}
    \text{myopic} &= \left[\frac{\delta_{ss'}^i}{p_{s'}^*} - \sum_{\tilde{s}'\in\mathcal{S}} \omega^*_{\tilde{s}'} \frac{\delta_{s\tilde{s}'}^i}{p_{\tilde{s}'}^*}\right] - \frac{\hat{\Lambda}(X,s,s')}{\Lambda^*(X,s,s')}  \\
    \text{hedging} &= \left[\frac{\hat{v}_i(X^*,s')}{v^*(X,s)}-\sum_{\tilde{s} \in \mathcal{S}}\omega^*_{\tilde{s}} 
    \frac{\hat{v}_i(X^*,\tilde{s})}{v^*(X,s)}\right]+
     \sum_{\tilde{s} \in \mathcal{S}}\omega^*_{\tilde{s}} \frac{\hat{\Lambda}(X,s,\tilde{s})}{\Lambda^*(X,s,\tilde{s})}.
\end{align}
\end{enumerate}
\end{proposition}
\begin{proof}
See Online Appendix \ref{proof: oa_investor_perturbation}.
\end{proof}

Proposition \ref{oa_investor_perturbation} provides asymptotic closed-form solutions to the value function and policy functions. 
The net-worth multiplier $\hat{v}_i(X,s)$ is high when investor $i$ is relatively optimistic and state-prices are relatively low. 
The effect of beliefs can be seen by writing the term involving $\delta_{ss'}^i$ as follows:
\begin{equation}
    \sum_{s'\in \mathcal{S}} \omega^*_{s'} \frac{1}{1-\gamma} \frac{\delta_{ss'}^i}{p_{s'}^*} = \sum_{s'\in \mathcal{S}} p_{s'}^* \frac{1}{1-\gamma} \frac{x_{s'}^{1-\gamma}}{\mathbb{E}^*[x_{s'}^{1-\gamma}]} \frac{\delta_{ss'}^i}{p_{s'}^*}  =  Cov^*\left(\frac{1}{1-\gamma} \frac{x_{s'}^{1-\gamma}}{\mathbb{E}^*[x_{s'}^{1-\gamma}]},\frac{\delta_{ss'}^i}{p_{s'}^*} \right),
\end{equation}
using the fact that $\sum_{s'\in \mathcal{S}} p_{s'}^* \frac{\delta_{ss'}^i}{p_{s'}^*} = 0$. 
The covariance above will be positive when $\delta_{ss'}^i$ is on average positive when $x_{s'}$ is high, i.e. the covariance is increasing in how optimistic investor $i$ is. 

The term involving $\hat{\Lambda}(X,s,s')$ captures the effect of changes in the SDF on the portfolio return that can be achieved by the household:\small
\begin{equation}
    1 = \mathbb{E}_s\left[ \Lambda(X,s,s';\epsilon) R_{i,n}(X,s,s';\epsilon) \right] \Rightarrow
    \sum_{s'\in\mathcal{S}} \omega_{s'}^* \frac{\hat{R}_{i,n}(X,s,s')}{R^*_{i,n}(X,s,s')} = - \sum_{s'\in\mathcal{S}} \omega_{s'}^* \frac{\hat{\Lambda}(X,s,s')}{\Lambda^*(X,s,s')}.
\end{equation}

Hence, if $\sum_{s'\in\mathcal{S}} \omega_{s'}^* \frac{\hat{\Lambda}(X,s,s')}{\Lambda^*(X,s,s')}$ is negative, then the household is able to achieve higher weighted portfolio returns in the $\epsilon > 0$ economy.

Given $\hat{v}_i(X,s)$, we can characterize the policy functions. The consumption-wealth ratio $\hat{c}_i(X,s)$ is proportional to $\hat{v}_i(X,s)$. If $\psi > 1$, such that the substitution effect on savings dominates,  households save a larger fraction of their wealth when average portfolio returns are high. 

As in the continuous-time model of \citet{merton1992continuous}, portfolio returns have two components: the \textit{myopic demand} and the \textit{hedging demand}. 
The myopic demand depends on current market conditions, while the hedging demand depends on future expected returns as captured by $\hat{v}_i(X^*,s')$.

We consider next the labor market outcomes and firms'  profits. 
\begin{proposition}[Hours, wages, and profits]\label{oa_labor_mkt} Suppose that $p_{ss'}^i = p_{s'}^* + \delta_{ss'}^i \epsilon$, $\alpha = \hat{\alpha} \epsilon$, and $\xi = \hat{\xi} \epsilon$. Suppose also that $\beta^* < 1$. Then,
\begin{enumerate}
    \item Wages:
\begin{equation}
    \hat{w}(E) = \hat{\xi}\left( \frac{\hat{\alpha} E}{\hat{\xi}} \right)^{\frac{\nu}{1+\nu}}.
\end{equation}

\item Hours:
\begin{equation}
    \hat{h}(E) = \left( \frac{\hat{\alpha} E}{\hat{\xi}} \right)^{\frac{1}{1+\nu}} \frac{\log \left( \frac{\hat{\alpha} E}{\hat{\xi}} \right)}{(1+\nu)^2} \hat{\alpha}.
\end{equation}

\item Profits:
\begin{equation}
    \hat{\pi}(X,s) = \left[ x_s \frac{\log \left( \hat{\alpha} E/   \hat{\xi} \right)}{1+\nu} -E \right] \hat{\alpha}.
\end{equation}
\end{enumerate}
\end{proposition}
\begin{proof}
See Online Appendix \ref{proof: oa_labor_mkt}
\end{proof}

We consider next the behavior of the price of the surplus claim and the riskless asset.
\begin{proposition}[Asset Prices]\label{oa_asset_prices} Suppose that $p_{ss'}^i = p_{s'}^* + \delta_{ss'}^i \epsilon$, $\alpha = \hat{\alpha} \epsilon$, and $\xi = \hat{\xi} \epsilon$. Suppose also that $\beta^* < 1$. Then,
\begin{enumerate}
    \item Price of surplus claim:
    \begin{equation}
        \frac{\hat{P}(X,s)}{P^*(X,s)}  = \left[\log (\hat{\alpha} E / \hat{\xi}) -  \frac{E}{x_s}\right] \frac{\hat{\alpha}}{1+\nu} + (\psi-1) \sum_{i=1}^I \eta_i \frac{\hat{v}_i(X,s)}{v^*(X,s)}.
    \end{equation}
    \item Return on the surplus claim:
    \begin{equation}
    \frac{\hat{R}_{p}(X,s,s')}{R_p^*(X,s,s')} = 
    \left[ \log \frac{E^*}{E} - \left( \frac{E^*}{x_{s'}} - \frac{E}{x_s} \right) \right]\frac{\hat{\alpha}}{1+\nu} + (\psi-1) \sum_{i=1}^I \eta_i \left[ \frac{\hat{v}_i(X^*,s')}{v^*(X^*,s')}- \frac{1}{\beta^*} \frac{\hat{v}_i(X,s)}{v^*(X,s)}\right].
\end{equation}
\item Risk-free rate:
    \begin{equation}
        \frac{\hat{R}_b(X,s)}{R_b^*(X,s)} = - \sum_{s'\in\mathcal{S}} \frac{p_{s'} x_{s'}^{-\gamma}}{\mathbb{E}^*[x_{s'}^{-\gamma}]} \frac{\hat{\Lambda}(X,s,s')}{\Lambda^*(X,s,s')}.
    \end{equation}
\item Conditional risk premium:
\begin{equation}
    \frac{\widehat{\overline{R}}_E(X,s)}{\overline{R}_E^*(X,s)} = \sum_{s'\in \mathcal{S}} \frac{p_{s'}^* x_{s'}}{\mathbb{E}^*[x_{s'}]} \frac{\hat{R}_E(X,s,s')}{R_E^*(X,s,s')} -\frac{\hat{R}_b(X,s)}{R^*_b(X,s)},
\end{equation}
where $\overline{R}_E(X,s;\epsilon) = \mathbb{E}^*\left[ \frac{R_E(X,s,s';\epsilon)}{R_b(X,s;\epsilon)}\right] $.
\end{enumerate}
\end{proposition}
\begin{proof}
See Online Appendix \ref{proof: oa_asset_prices}
\end{proof}

The next proposition provides the law of motion of the aggregate state variables.
\begin{proposition}[Aggregate state variables.]\label{oa_aggregate_state} Suppose that $p_{ss'}^i = p_{s'}^* + \delta_{ss'}^i \epsilon$, $\alpha = \hat{\alpha} \epsilon$, and $\xi = \hat{\xi} \epsilon$. Suppose also that $\beta^* < 1$. Then,
\begin{enumerate}
    \item Wealth distribution:
    \begin{equation}
            \frac{\hat{\eta}_i'(X,s,s')}{\eta_i} =   \frac{\hat{R}_{i,n}(X,s,s')}{R^*_{i,n}(X,s,s')}-\sum_{j=1}^I \eta_i \frac{\hat{R}_{j,n}(X,s,s')}{R^*_{j,n}(X,s,s')} - ((\beta^*)^{-1}-1)\left(\frac{\hat{c}_i(X,s)}{c_i^*(X,s)}  - \sum_{j=1}^I \eta_j \frac{\hat{c}_j(X,s)}{c_j^*(X,s)} \right) + \kappa \frac{\mu_i - \eta_i}{\eta_i}.
    \end{equation}
    \item Risk-neutral expectation of productivity growth:
    \begin{equation}
    \frac{\hat{E}'(X,s)}{E^*} = \frac{\hat{R}_b(X,s)}{R_b^*(X,s)}  + \sum_{s'\in \mathcal{S}} \omega_{s'}^*\frac{\hat{\Lambda}(X,s,s')}{\Lambda^*(X,s,s')} .
\end{equation}
\end{enumerate}

\begin{proof}
See Online Appendix \ref{proof: oa_aggregate_state}.
\end{proof}
\end{proposition}

Propositions \ref{oa_investor_perturbation} to \ref{oa_aggregate_state} characterize the behavior of all equilibrium objects given the economy's SDF $\hat{\Lambda}(X,s,s')$. 
The next proposition provides an expression for $\Lambda(X,s,s')$ in terms of the primitives of the economy.

\begin{proposition}[The economy's SDF]\label{oa_SDF} Suppose that $p_{ss'}^i = p_{s'}^* + \delta_{ss'}^i \epsilon$, $\alpha = \hat{\alpha} \epsilon$, and $\xi = \hat{\xi} \epsilon$. Suppose also that $\beta^* < 1$. Then, the economy's SDF is given by\small
\begin{equation}
    \frac{\hat{\Lambda}(X,s,s')}{\Lambda^*(X,s,s')} = \gamma b^{\Lambda}(X,s,s') -(\gamma - \psi^{-1}) \left[  \omega^* b^{\Lambda}(X,s) -  \beta^* \omega^* \cdot b^\Lambda(X^*,s') + \beta^* \sum_{\tilde{s}\in \mathcal{S}} \omega_{\tilde{s}} ( \omega^* \cdot b^{\Lambda}(X^*,\tilde{s}) )
    \right],
\end{equation}\normalsize
where \footnotesize
\begin{equation}
    b^\Lambda(X,s,s')  = \frac{1}{\gamma}\frac{\delta_{ss'}(X)}{p_{s'}^*}- \frac{\psi - \gamma^{-1}}{\gamma-1} \left[\omega^* \cdot \delta_s(X) - \beta^* \omega^* \cdot \delta_{s'}(X)+\beta^* \sum_{\tilde{s}} \omega_{\tilde{s}}^* (\omega^* \cdot \delta_{\tilde{s}}(X) ) \right]
- \left[ \log\frac{E^*}{E} - \left( \frac{E^*}{x_{s'}} - \frac{E}{x_s} \right) \right]\frac{\hat{\alpha}}{1+\nu}.
\end{equation}\normalsize
\end{proposition}
\begin{proof}
See Online Appendix \ref{proof: oa_SDF}.
\end{proof}

A particularly simple special case is given by the case of CRRA preferences, i.e. $\gamma = \psi^{-1}$.
\begin{corollary} Suppose $\gamma = \psi^{-1}.$ Then,
\begin{enumerate}
    \item SDF:
    \begin{equation}
    \frac{\hat{\Lambda}(X,s,s')}{\Lambda^*(X,s,s')} = \frac{\delta_{ss'}(X)}{p_{s'}^*} - \gamma  \left[ \log\frac{E^*}{E} - \left( \frac{E^*}{x_{s'}} - \frac{E}{x_s} \right) \right]\frac{\hat{\alpha}}{1+\nu}.
\end{equation}
%     \item Net-worth multiplier:
%     \begin{equation}
%     \frac{\hat{v}_i(X,s)}{v_i^*(X,s)}  =  \beta^*  \sum_{s'\in \mathcal{S}} \omega^*_{s'}\left[\frac{1}{1-\gamma} \frac{\delta_{ss'}^i}{p_{s'}^*} -\frac{\hat{\Lambda}(X,s,s')}{\Lambda^*(X,s,s')} \right] +\beta^* \overline{v},
% \end{equation}
% where $\omega_{s'}^* \equiv\frac{p_{s'}^* x_{s'}^{1-\gamma}}{\sum_{\tilde{s} \in \mathcal{S}} p_{\tilde{s}}^* x_{\tilde{s}}^{1-\gamma}}$,$X^* = (E^*, \{\eta_i\}_{i=1}^I)$, and
% \begin{equation}
%     \overline{v} \equiv \frac{ \beta^* }{1-\beta^*}\sum_{\tilde{s}\in \mathcal{S}}  \omega^*_{\tilde{s}}  \sum_{s'\in \mathcal{S}} \omega^*_{s'}\left[\frac{ 1}{1-\gamma} \frac{\delta_{\tilde{s}s'}^i}{p_{s'}^*} -\frac{\hat{\Lambda}(X^*,\tilde{s},s')}{\Lambda^*(X^*,\tilde{s},s')}  \right].
% \end{equation}
\end{enumerate}
\end{corollary}


\subsection{Special cases}

Consider the special case where $\delta_{ss'}^i = 0$ for all $i \in \mathcal{I}$ and $s, s' \in \mathcal{S}$. In this case, investors still have common iid beliefs, but returns will not be iid due to the fact that labor is chosen one period in advance. 

In this case, the economy's is given by
\begin{equation}
    \frac{\hat{\Lambda}(X,s,s')}{\Lambda^*(X,s,s')} =\psi^{-1} \left[ \log \frac{E}{E^*} +\frac{E^*}{x_{s'}}  - \frac{E}{x_s}\right]  \frac{\hat{\alpha}}{1+\nu} - (\gamma - \psi^{-1}) (1-\beta^*) \sum_{\tilde{s}'}\omega^*_{\tilde{s}'} \left(  \frac{E^*}{x_{\tilde{s}'}} - \frac{E^*}{x_{s'}} \right) \frac{\hat{\alpha}}{1+\nu}.
\end{equation}

The interest rate is given by
\begin{equation}
        \frac{\hat{R}_b(X,s)}{R_b^*(X,s)} = - \sum_{s'\in\mathcal{S}} \frac{p_{s'} x_{s'}^{-\gamma}}{\mathbb{E}^*[x_{s'}^{-\gamma}]} \frac{\hat{\Lambda}(X,s,s')}{\Lambda^*(X,s,s')}.
    \end{equation}

\subsection{Conditional moments}


Consider the conditional risk premium
\begin{equation}
    R^e_p(X,s;\epsilon) \equiv \sum_{s'\in\mathcal{S}} p_{s'}^* \frac{R_p(X,s,s';\epsilon)}{R_b(X,s;\epsilon)}.
\end{equation}

Expanding the expression above in $\epsilon$, we obtain
\begin{equation}
    \frac{\hat{R}_p^e(X,s)}{R_p^{e,*}} = \sum_{s'\in\mathcal{S}} \frac{p_{s'}^* R_p^*(X,s,s')}{\sum_{\tilde{s}'\in \mathcal{S}} p_{\tilde{s}'}^* R_p^*(X,s,\tilde{s}')} \left[ \frac{\hat{R}_p(X,s,s')}{R_p^*(X,s,s')} - \frac{\hat{R}_b(X,s)}{R_b^*(X,s)}\right] = 
    \sum_{s'\in \mathcal{S}} \frac{p_{s'}x_{s'}}{\mathbb{E}^*[x_{s'}]} \frac{\hat{R}_p(X,s,s')}{R_p^*(X,s)} - \frac{\hat{R}_b(X,s,s')}{R_b^*(X,s)}.
\end{equation}

The conditional volatility of excess returns is given by
\begin{equation}
    \sigma_{p}(X,s;\epsilon) \equiv \left[ \sum_{s'\in\mathcal{S}} p_{s'} \left( R_p^e(X,s,s';\epsilon)- R_p^e(X,s;\epsilon)\right)^2 \right]^{\frac{1}{2}}.
\end{equation}

Expanding the expression above in $\epsilon$, we obtain
\begin{equation}
    \frac{\hat{\sigma}_{p}(X,s)}{\sigma^*_p(X,s)} = \frac{1}{\sigma_p^*(X,s)^2}  \sum_{s'\in\mathcal{S}} p_{s'} \left( R_p^{e,*}(X,s,s') - R_p^{e,*}(X,s)\right)  \left( \hat{R}_p^{e}(X,s,s') - \hat{R}_p^{e}(X,s)\right),
\end{equation}
where $\hat{R}_p^{e}(X,s,s') = \hat{R}_p^{e,*}(X,s,s') \left( \frac{\hat{R}_p(X,s,s')}{R^*_p(X,s,s')} - \frac{\hat{R}_b(X,s)}{R_b^*(X,s)} \right) $.


\subsection{Stock prices in the log economy}

Suppose $\psi = \gamma = 1$ and that investors have homogeneous iid beliefs, $\delta_{ss'}^i = 0$. The stock price satisfies the following recursion:
\begin{equation}
    Q(X,s) = \sum_{s'\in \mathcal{S}} p_{s s'} \Lambda(X,s, s') \left[ \pi(\mathcal{L}'(X,s), s') + x_{s'} Q(X', s') \right],
\end{equation}
where $\pi(E, s) =  \left( \frac{\alpha E}{\xi} \right)^{\frac{\alpha}{1+\nu - \alpha}} x_s \left[ 1- \frac{\alpha}{x_s} E \right] $ and $p_{ss'} \Lambda (X,s,s') = p_{ss'}(X) \frac{\beta}{x_{s'}} \left( \frac{E}{\mathcal{L}'(X,s)}  \right)^{\frac{\alpha}{1+\nu - \alpha}} \frac{1+\nu - \alpha \frac{E}{x_s}}{1+\nu - \alpha \frac{\mathcal{L}'(X,s)}{x_{s'}}}$.

Define the price-dividend ratio $q(X,s) = x_s \frac{Q(X,s)}{\pi(X,s)}$. The price-dividend ratio satisfies the recursion:
\begin{equation}
    q(X,s) = \sum_{s'\in \mathcal{S}} p_{s s'} \Lambda(X,s, s') \left[ \pi(\mathcal{L}'(X,s), s') + x_{s'} Q(X', s') \right],
\end{equation}

We can then write the expression above as follows:
\begin{equation}
    q(X,s) =  \beta \sum_{s'\in \mathcal{S}} p_{s s'}(X)  \frac{1+\nu - \alpha \frac{E}{x_s}}{1+\nu - \alpha \frac{\mathcal{L}'(X,s)}{x_{s'}}} \frac{1 - \alpha \frac{\mathcal{L}'(X,s)}{x_{s'}}}{1 - \alpha \frac{E}{x_s}}\left[ 1+ q(X', s') \right].
\end{equation}

Let's assume that $\nu = \overline{\nu} \epsilon$. We can then write $q(X,s; \epsilon)$ as follows
\begin{equation}
    q(X,s;\epsilon) = \sum_{k=0}^\infty q_k(X,s) \epsilon^k. 
\end{equation}

Define $g(X,s,s';\epsilon) \equiv \frac{1+\nu - \alpha \frac{E}{x_s}}{1+\nu - \alpha \frac{\mathcal{L}'(X,s)}{x_{s'}}} \frac{1 - \alpha \frac{\mathcal{L}'(X,s)}{x_{s'}}}{1 - \alpha \frac{E}{x_s}}$. We can expand $g(X,s,s')$ as follows
\begin{equation}
    g(X,s,s';\epsilon) = \sum_{k=0}^\infty g_{k}(X,s,s') \epsilon^k, 
\end{equation}
where  $g_0(X,s,s') =1 $ and, for $k > 0$, we obtain
\begin{equation}
    g_k(X,s,s') =  \frac{\alpha\overline{\nu}^k \left(\frac{\mathcal{L}'(X,s)}{x_{s'}}- \frac{E}{x_s}\right) }{\left(1 - \alpha\frac{E}{x_s}\right) \left( \alpha\frac{\mathcal{L}'(X,s')}{x_{s'}}-1\right)^k } 
\end{equation}

This gives the following recursion for $q_{k}(X,s)$:
\begin{equation}
    q_k(X,s) = \beta \sum_{s'\in \mathcal{S}} p_{s s'}(X)  \left[ g_k(X,s)+ \sum_{j=0}^k g_j(X,s) q_{k-j}(X', s') \right].
\end{equation}

Under our assumptions, the risk-neutral expectation of productivity growth is constant $\mathcal{L}'(X,s) = \overline{E}$.  In this case, we can write the recursion
\begin{equation*}
    Q(\overline{X}, s) = \beta \left( 1 - \frac{\alpha}{1+\nu} \frac{\overline{E} }{x_s}\right) \sum_{s' \in \mathcal{S}} \frac{p^*_{s'} x_{s'}}{x_{s'} - \frac{\alpha}{1+\nu} \overline{E}} \left[ Q(\overline{X}, s')  + \left( \frac{\alpha \overline{E}}{\xi}  \right)^\frac{\alpha}{1+\nu -\alpha} \left( 1 - \alpha  \frac{\overline{E}}{x_{s'}} \right)\right]
\end{equation*}

Let $\tilde{Q}(X,s) \equiv \frac{Q(x,s)}{1 - \frac{\alpha}{1+\nu} \frac{\mathcal{L}'(X,s)}{x_s}}$,
$\tilde{Q}(X) \equiv [\tilde{Q}(X,1), \ldots, \tilde{Q}(X,N)]' $, and $b^Q \equiv \beta \left( \frac{\alpha \overline{E}}{\xi}  \right)^\frac{\alpha}{1+\nu -\alpha} \sum_{s'\in \mathcal{S}} p_{s'}^* \frac{x_{s'} - \alpha \overline{E}}{x_{s'} - \frac{\alpha}{1+\nu}\overline{E} } $.
Then, we can write 
\begin{equation}
    \left[ I - \beta \mathbb{1}_N (p^*)' \right] \tilde{Q}(\overline{X}) = b^Q \mathbb{1}_N,
\end{equation}

Inverting the matrix above, we obtain
\begin{equation}
    Q(\overline{X},s) = \left[\frac{\beta}{1-\beta} \left( \frac{\alpha \overline{E}}{\xi}  \right)^\frac{\alpha}{1+\nu -\alpha}   \sum_{s'\in \mathcal{S}} p_{s'}^* \frac{x_{s'} - \alpha \overline{E}}{x_{s'} - \frac{\alpha}{1+\nu}\overline{E}} \right]\left( 1 - \frac{\alpha}{1+\nu} \frac{\overline{E} }{x_s}\right). 
\end{equation}

The price-dividend ratio is given by
\begin{equation}
    \frac{x_s Q(\overline{X},s)}{\pi(\overline{X},s)} =  \frac{\beta}{1-\beta} \sum_{s'\in \mathcal{S}} p_{s'}^* \frac{x_{s'} - \alpha \overline{E}}{x_{s'} - \frac{\alpha}{1+\nu}\overline{E}} \frac{x_s - \frac{\alpha}{1+\nu} \overline{E}}{x_s - \alpha\overline{E}}.
\end{equation}

Equity returns are given by
\begin{equation}
    R_E(X,s, s') = \frac{\pi(\overline{E}, s') + x' Q(\overline{X}, s')}{Q(\overline{X}, s)} = \frac{x_{s'}}{\beta}\left[  \frac{1-\beta}{\sum_{s'\in \mathcal{S}} p_{s'}^* \frac{x_{s'} - \alpha \overline{E}}{x_{s'} - \frac{\alpha}{1+\nu}\overline{E}} } \frac{1 - \alpha \frac{\overline{E}}{x_{s'}}}{1 - \frac{\alpha}{1+\nu} \frac{\overline{E}}{x_{s}} } + \beta \frac{1 - \frac{\alpha}{1+\nu} \frac{\overline{E}}{x_{s'}}}{1 - \frac{\alpha}{1+\nu} \frac{\overline{E}}{x_{s}} } \right].
\end{equation}

Excess returns are given by
\begin{equation}
    R_E^e(X,s,s') = a_E x_{s'} + b_E,
\end{equation}
where
\begin{align}
    a_E &\equiv  \frac{1}{\left(1-\frac{\alpha}{1+\nu}\right) \overline{E} } \left( \frac{1-\beta}{\sum_{s'\in \mathcal{S}} p_{s'}^* \frac{x_{s'}  - \alpha \overline{E}}{x_{s'} - \frac{\alpha}{1+\nu}\overline{E}} } + \beta \right)\\
    b_E &\equiv 
\end{align}

The conditional risk premium is given by
\begin{equation}
    R_E(\overline{X},s) = \frac{1}{\left(1-\frac{\alpha}{1+\nu}\right) \overline{E}}\left[ \left( \frac{1-\beta}{\sum_{s'\in \mathcal{S}} p_{s'}^* \frac{x_{s'}  - \alpha \overline{E}}{x_{s'} - \frac{\alpha}{1+\nu}\overline{E}} } + \beta \right) \mathbb{E}[x_{s'}]-\left( \frac{(1-\beta)(1+\nu)}{\sum_{s'\in \mathcal{S}} p_{s'}^* \frac{x_{s'}  - \alpha \overline{E}}{x_{s'} - \frac{\alpha}{1+\nu}\overline{E}} } + \beta \right)  \frac{\alpha}{1+\nu} \overline{E}\right]-1.
\end{equation}

We can write the expression above as follows:
\begin{equation}
    R_E(\overline{X},s) = \frac{1}{\left(1-\frac{\alpha}{1+\nu}\right) \overline{E}}\left[ \left( \frac{1-\beta}{1 - \nu \sum_{s'\in \mathcal{S}} \frac{p_{s'}^* \frac{\alpha}{1+\nu} \overline{E}}{x_{s'} - \frac{\alpha}{1+\nu}\overline{E}} } + \beta \right) \mathbb{E}[x_{s'}]-\left( \frac{(1-\beta)(1+\nu)}{1 - \nu \sum_{s'\in \mathcal{S}} \frac{p_{s'}^* \frac{\alpha}{1+\nu} \overline{E}}{x_{s'} - \frac{\alpha}{1+\nu}\overline{E}}} + \beta \right)  \frac{\alpha}{1+\nu} \overline{E}\right]-1.
\end{equation}



\subsection{Quantitative implications}

Let $z_t$ denote demeaned log productivity growth, which we assume follows an AR(1) process:
\begin{equation}
    z_{t+1} = \rho z_t + \sigma \sqrt{1-\rho^2}\epsilon_{t+1},
\end{equation}
where $\epsilon_{t+1}$ follows a standard normal distribution and it is serially uncorrelated. In levels, the (gross) productivity growth is given by $x_t = e^{\mu+z_t} $, where $\mu$ denotes average productivity growth. 

We discretize the process above following the method of \citet{rouwenhurst1995asset}. Let $\hat{z}_t$ denote the discretized variable taking values in the equally-spaced grid  $\{z_{1}, \ldots, z_N \}$, where $z_i = -\overline{\psi}  + \frac{2 \overline{\psi}}{N-1} (i-1) $, so $z_1 = -\overline{\psi}$ and $z_N =  \overline{\psi}$. We set $\overline{\psi} \equiv \sigma \sqrt{N-1}$, so we match the unconditional variance. 

\subsection{A more general process for productivity growth}\label{sec:discretization}

% comment move paragraph entirely to appendix. It is not essential. 
\paragraph{Discretization.}
The evolution of $\hat{x}_t$, under subjective beliefs, can be written in a convenient matrix form:
\begin{equation}\label{eq: MS_VAR}
    \left[ 
    \begin{matrix}
        \hat{x}_{t+1}\\
        \hat{z}_{i,t+1}
    \end{matrix}
    \right] = \left[ \begin{matrix}
        w_{t+1}\\
        0
    \end{matrix}\right]  + \left[ \begin{matrix}
        \theta_i & 1 \\
        0 & 0
    \end{matrix} \right]
    \left[ 
    \begin{matrix}
        \hat{x}_{t}\\
        \hat{z}_{i,t}
    \end{matrix}
    \right] + \left[ \begin{matrix}
        \sigma_{i,u} & 0\\
        0 & \sigma_v
    \end{matrix} \right] \left[ 
    \begin{matrix}
        u_{i,t+1}\\
        v_{t+1}
    \end{matrix}
    \right],
\end{equation}
where $\hat{z}_{i,t} \equiv \mathbb{E}_{i,t}[\hat{x}_{t+1}] - \theta_i \hat{x}_{t}$. We recover objective beliefs in the special case $\theta_i = \sigma_v = 0$.
As $w_{t+1}$ follows a Markov chain, the process above corresponds to a Markov-switching vector autoregression (MS-VAR), with state-dependent conditional means. To discretize the system above, we adapt the methods of \citet{gospodinov2014moment}, who extended the \citet{rouwenhurst1995asset} method to VARs, and \citet{liu2017discretization}, who proposed a discretization of univariate Markov-Switching models. 
The discretization provides a state space with dimension $S$  for $x_t$, so $x_t \in \mathcal{X}= \{x^1, x^2, \ldots, x^S\}$,  and transition probabilities $\{p_{ss'}^i\}$, for $s, s' \in \mathcal{S} = \{1,2, \ldots, S\}$, that approximate the MS-VAR \eqref{eq: MS_VAR}. Notice that our discretization implies that the grid $\mathcal{X}$ is the same for all investors, so they agree on the state $s$, but they disagree on the transition probabilities $p_{ss'}^i$.



Let $\hat{x}_t \equiv \log x_t - \mu $  denote the demeaned log productivity growth. We assume that investor $i$ believes the process for $\hat{x}_t$ is given by
\begin{align}
    \hat{x}_{t+1} &= \mathbb{E}_{i,t}[\hat{x}_{t+1}] + \sigma_{i,u} u_{i,t+1}\\
    \mathbb{E}_{i,t}[\hat{x}_{t+1}] &= \theta_i \hat{x}_t + \sigma_{v} v_{i,t},
\end{align}
$v_{i,t} = \overline{v}_t + \Tilde{v}_{i,t}$, 
where $u_{i,t}$ and $v_t$ are mutually independent, serially uncorrelated, standard normal random variables. 
Notice that $u_{i,t+1}$ represents the period $t+1$ innovation according to investor $i$ and $v_t$ represents an \textit{expectation shock}. We assume that this expectation shock is common across investors, so heterogeneity comes only from $\theta_i$. 

The presence of this expectation shock is important to quantitatively match the volatility of expectations in the data. 
To see the role of $v_t$, notice that the unconditional variance of $\hat{x}_{t+1}$ and $\mathbb{E}_t[\hat{x}_{t+1}]$ are given by
\begin{equation}
    Var[\hat{x}_{t+1}] = \frac{\sigma_{i,u}^2 + \sigma_v^2}{1-\theta_i^2}, \qquad \qquad Var[\mathbb{E}_t[\hat{x}_{t+1}]] =  \frac{\theta_i^2 \sigma_{i,u}^2 + \sigma_v^2}{1-\theta_i^2}.
\end{equation}

The fraction of total variance explained by movements in expectations is given by
\begin{equation}
    \frac{Var[\mathbb{E}_t[\hat{x}_{t+1}]]}{Var[\hat{x}_{t+1}]} = \theta_i^2 + (1-\theta_i^2) \frac{\sigma_v^2}{\sigma_{i,u}^2 + \sigma_v^2}.
\end{equation}
Hence, by adjusting $\sigma_v$, it is possible to obtain any value in the interval $[\theta_i^2, 1)$ for the fraction of variance explained by movements in expectations. In the special case $\sigma_v = 0$, we obtain an AR(1) process for $\hat{x}_{t+1}$, which achieves the lower bound of this interval. 



\paragraph{Discretization.} We discretize the process above using the generalization of the method of \citet{rouwenhurst1995asset} proposed by \citet{gospodinov2014moment}. 
The method consists of mixing the distribution for independent AR(1) processes to approximate the distribution of a VAR(1) with uncorrelated shocks. 
Define $\hat{z}_{t} \equiv \mathbb{E}_{i,t}[\hat{x}_{i,t}] - \theta_i \hat{x}_t$, so we can write the system above in matrix form:
\begin{equation}
    \left[ \begin{matrix}
        \hat{x}_{t+1}\\
        \hat{z}_{t+1}
    \end{matrix}\right] = \left[ \begin{matrix}
        \theta & 1\\
        0 & 0
    \end{matrix}\right]
    \left[ \begin{matrix}
    \hat{x}_t\\
    \hat{z}_{t}
    \end{matrix}
    \right]+\left[ \begin{matrix}
        \sigma_u & 0\\
        0 & \sigma_v
    \end{matrix}\right]
    \left[ \begin{matrix}
     u_{t+1}\\
     v_{t+1}
    \end{matrix}
    \right],
\end{equation}
where we dropped the dependence on the investor $i$ to ease notation.
Given this representation, we can construct the discrete approximation following the three steps described below.


\textit{Step 1: grid construction.}
We construct the grids for $\hat{x}$ and $\hat{z}$ as in \citet{rouwenhurst1995asset}. 
Let $\hat{x}(N_x, \sigma_x) = \{ \overline{x}^1, \overline{x}^2, \ldots, \overline{x}^{N_x} \}$ denote the grid for $\hat{x}$, where 
\begin{equation}
    \hat{x}^i = - \psi_x(N_x, \sigma_x) + 2 \psi_x(N_x, \sigma_x) \frac{i-1}{N_x-1}, 
\end{equation}
$\psi_x(N_x, \sigma_x) \equiv \sigma_x \sqrt{N_x -1}$, and $\sigma_x$ denotes the unconditional standard-deviation for $\hat{x}_t$. Notice that the grid is equally spaced, $\overline{x}^1 = - \psi_x(N_x, \sigma_x)$, and $\overline{x}^{N_x} = \psi_x(N_x, \sigma_x)$. 
The grid for $\hat{z}$ is constructed analogously. 


\textit{Step 2: transition matrix for independent AR(1).} 
Let $\Pi(N, \rho, \sigma)$ denote the $N\times N$ transition matrix for the \citet{rouwenhurst1995asset} approximation of an AR(1) process with autocorrelation $\rho$ and unconditional variance $\sigma^2$. We denote the $k$-th row of this matrix by 
\begin{equation}
    \pi_k(N, \rho, \sigma) = \{\pi_{k,1}(N, \rho, \sigma), \pi_{k,2}(N, \rho, \sigma), \ldots, \pi_{k,N}(N, \rho, \sigma)\},
\end{equation}
where $\pi_{k,l}(N, \rho, \sigma)$ is the probability of transitioning from state $k$ to state $l$. In the special case where $\rho = 0$, the transition probability is independent of the current state, so we can write $\pi_{k,j}(N,0,\sigma) = \overline{\pi}_j(N,\sigma)$.


\textit{Step 3: Markov chain  construction.} Given $N_x$ points in the grid for $\hat{x}$ and $N_{z}$ points in the grid for $\hat{z}_t$, we have a total of $S = N_x \times N_z$ states. Denote the state space by $\mathcal{S} = \{1,2, \ldots, S\}$. Let $s = i + (k-1) \times N_x$ and $s' = j + (l-1) \times N_x$, where $i, j \in \{1, \ldots, N_x\}$ and $k, l \in \{1, \ldots, N_z\}$. 
Denote the probability of $\hat{x}_{t+1} = \overline{x}^{j}$ given state $s$ by $p_s^x(j)$ and the probability of $\hat{z}_{t+1} = \overline{z}^{l}$ given state $s$ by $p_{s}^z(l)$. 
As $\hat{x}_{t+1}$ and $\hat{z}_{t+1}$ are conditionally independent, the probability of switching from state $s$ to state $s'$ is given by 
\begin{equation}
    p_{ss'} = p_s^x(j) \times p_s^{z}(l).
\end{equation}

As $z$ is serially uncorrelated, we have that $p_{s}^z(l) = \overline{\pi}_l(N_z, \sigma_z)$.
The transition probability for $\hat{x}_t$ will be obtained by appropriately mixing the distribution of an AR(1) process with autocorrelation $\rho_{x} \equiv \sqrt{1 - \frac{\sigma_u^2}{\sigma_x^2}}$ and unconditional variance $\sigma_x^2$.

Let $\mu_x(s) \equiv \theta \hat{x}^{i} + \hat{z}^{k}$ denote the conditional expectation of $\hat{x}$ at state $s$. Suppose first that $\mu_x(s) \in [ \rho_x \overline{x}^1, \rho_x \overline{x}^{N_x} ]$. Define the probability of $\hat{x}_{t+1} = \overline{x}^{j}$ given state $s$ as follows:
\begin{equation}
p_s^x(j)  = \lambda(\rho_x) \pi_{\iota,j}(N_x, \rho_x, \sigma_x) + (1-\lambda(\rho_x)) \pi_{\iota+1,j}(N_x, \rho_x, \sigma_x),
\end{equation}
where $\iota$ is such that  $\rho_{x} \overline{x}^\iota \leq \mu_x(s) \leq \rho_x \overline{x}^{\iota+1}$ and $\lambda(\rho_x) \equiv \frac{\rho_{x} \overline{x}^{\iota+1} - \mu_x(s)}{\rho_{x} \overline{x}^{\iota+1}-\rho_{x} \overline{x}^\iota}$.

This choice of $\lambda(\rho_x)$ implies that we match the conditional moments:
\begin{equation}
     \sum_{j=1}^{N_x} p_{s}^x(j)\overline{x}^{j}  = \lambda(\rho_x) \rho_x \overline{x}^{\iota} + (1-\lambda(\rho_x)) \rho_x \overline{x}^{\iota+1} = \mu_{x}(s).
\end{equation}

The conditional second moment is given by
\begin{equation}
     \sum_{j=1}^{N_x} p_s^x(j) (\overline{x}^{j})^2 = \sigma_x^2 (1-\rho_x^2) + \rho_x^2 \left[\lambda(\rho_x) (\overline{y}^{\iota})^2 + (1-\lambda(\rho_x)) (\overline{y}^{\iota+1})^2\right].
\end{equation}

Denote the conditional variance of the discrete process by $\Tilde{\sigma}_u^2$, which is given by 
\begin{align}
    \Tilde{\sigma}_u^2 &=  \sigma_x^2 (1-\rho_x^2) + \rho_x^2 \left[\lambda(\rho_x) (\overline{y}^{\iota})^2 + (1-\lambda(\rho_x)) (\overline{y}^{\iota+1})^2\right]  - \rho_x^2 \left[\lambda(\rho_x) \overline{y}^{\iota} + (1-\lambda(\rho_x)) \overline{y}^{\iota+1}\right]^2 \nonumber\\
    &= \sigma_x^2 (1-\rho_x^2) + \rho_x^2 \lambda(\rho_x) (1-\lambda(\rho_x)) ( \overline{x}^{\iota+1} - \overline{x}^{\iota} )^2\nonumber\\
    &= \sigma_x^2 (1-\rho_x^2) + \sigma_x^2 \rho_x^2 \frac{4\lambda(\rho_x) (1-\lambda(\rho_x)) }{N_x-1},
\end{align}
using the fact that $(\overline{x}^{i+1} - \overline{x}^{i})^2 = \frac{4 \sigma_x^2}{N-1}$.
As $N_x \rightarrow \infty$, the second term on the right converges to zero and $\Tilde{\sigma}_u = \sigma_x^2 (1-\rho_x^2) = \sigma_u^2$, given our choice of $\rho_x$. If $\mu_{x}(s)/\rho_x$ does not belongs to the grid of $\hat{x}$, then the discretization matches the conditional mean of $\hat{x}$, but it overstates the conditional variance.  

Suppose now that $\mu_x(s) \notin [\rho_x \overline{x}^1, \rho_x \overline{x}^{N_x}]$. In this case, we set $p_s^x(j) = \pi_{1,j}(N_x, \rho_x, \sigma_x)$ if $\mu_{x}(s) < \rho_x \overline{x}^1$ and $p_s^x(j) = \pi_{N_x,j}(N_x, \rho_x   , \sigma_x)$ if $\mu_x(s) > \rho_x \overline{x}^{N_x}$. In both cases, the conditional variance is matched exactly and the conditional mean achieves the value closest to $\mu_{x}(s)$ given the grid points. 

\paragraph{A different representation.} An equivalent representation of the system is given by
\begin{align}
    \hat{x}_{t+1} &= z_{t} + \sigma_u u_{t+1}\\
    z_{t+1} &= \theta_i z_{t} + \theta_i \sigma_u u_{t+1}+ \sigma_{v} v_{t+1},
\end{align}
where $z_t \equiv \mathbb{E}_{i,t}[\hat{x}_{t+1}]$. Hence, expected growth follows an AR(1) process and it is exposed to both expectation shocks, $v_{t+1}$, and shocks to realized growth rates, $u_{t+1}$. Notice that we cannot independently choose the persistence of expectations and the correlation between $z_{t+1}$ and $\hat{x}_{t+1}$.

The impact of $v_t$ in expected future growth is
\begin{equation}
    \frac{\partial \mathbb{E}_t[\hat{x}_{t+k}]}{\partial v_t} = \sigma_v \theta_i^{k-1},
\end{equation}
for $k \geq 1$.

\subsection{A process with richer heterogeneity}

Under the objective measure, log productivity follows a Markov-Switching process:
\begin{equation}
    \log(x_{t+1}) = \mu_{t+1} + \theta ( \log (x_{t}) - \mu_t) + u_{t+1},
\end{equation}
where $u_{t+1} \sim \mathcal{N}(0, \sigma_u^2)$ and 
$\mu_{t+1}$ follows a two-state Markov chain, that is, $\mu_{t+1}\in \{\mu^1, \mu^2\}$ and $Pr(\mu_{t+1} = \mu^j|\mu_t = \mu^i) = p^\mu_{ij}$ for $i,j \in \{1,2\}$. 
The different regimes enable us to capture the fact that productivity is subject to small fluctuations most of the time with occasional rare large shocks.

Under subjective beliefs, productivity follows the process
\begin{align}
    \log(x_{t+1})  &= \mu_{t+1} + \theta_i (\log (x_{t}) - \mu_t) + v_{i,t}+ u_{i,t+1},
\end{align}
where $u_{i,t+1} \sim \mathcal{N}(0, \sigma_{i,u}^2)$, $v_{i,t} = \rho \sigma_{i,v} \overline{v}_t + \sqrt{1-\rho^2} \sigma_{i,v}\hat{v}_{i,t} $, and $(\overline{v}_t,v_{i,t})$ are iid standard normal random variables. We assume that $(u_{i,t}, \hat{v}_{i,t}, \overline{v}_{t})$ are mutually independent.

Subjective beliefs differ from the objective one in two important dimensions. First, the persistent parameter $\theta_i$ may differ from the objective one $\theta$. 
Second, subjective beliefs are exposed to expectation shocks $v_{i,t}$. 
These expectations shocks are exposed to a common component $\overline{v}_t$ and an investor-specific component $v_{i,t}$.
Differences in $\theta_i$ capture the fact that investors differ on how they react to news, with some investors extrapolating and some investors under-reacting. 
The expectation shocks $v_{i,t}$ are important to capture the volatility of subjective expectations observed in the data. 

Define  $\hat{x}_t \equiv \log (x_{t}) - \mu_t$ and the vector $\hat{v}_t = [\hat{v}_{1,t}, \ldots, \hat{v}_{I,t}]'$. Investor $i$ believes that $[\hat{x}_{t}, \overline{v}_t, \hat{v}_t]$ follows the process:
\begin{equation}
    \left[
    \begin{matrix}
        \hat{x}_{t+1}\\
        \overline{v}_{t+1}\\
        \hat{v}_{t+1}
    \end{matrix}
    \right] = 
    \left[ 
    \begin{matrix}
    \theta_i & \rho \sigma_{i,v} & \sqrt{1-\rho^2} \sigma_{i,v}\textbf{e}_i'\\
    0 & 0 & \textbf{0}_{1\times I}\\
    \textbf{0}_{I\times 1} & \textbf{0}_{I\times 1} & \textbf{0}_{I\times I}
    \end{matrix}
    \right] \left[
    \begin{matrix}
        \hat{x}_{t}\\
        \overline{v}_t\\
        \hat{v}_{t}
    \end{matrix}
    \right]  +\left[ 
    \begin{matrix}
    u_{i,t+1}\\
    \overline{v}_{t+1}\\
    \hat{v}_{t+1}
    \end{matrix}
    \right] 
\end{equation}

Notice that the total variance and the variance of the conditional expectation are given by
\begin{equation}
    Var[\hat{x}_{t+1}] = \frac{\sigma_{i,u}^2 + \sigma_{i,v}^2}{1-\theta_i^2}, \qquad \qquad Var[\mathbb{E}_t[\hat{x}_{t+1}]] =  \frac{\theta_i^2 \sigma_{i,u}^2 + \sigma_{i,v}^2}{1-\theta_i^2}.
\end{equation}

The fraction of total variance explained by movements in expectations is given by
\begin{equation}
    \frac{Var[\mathbb{E}_t[\hat{x}_{t+1}]]}{Var[\hat{x}_{t+1}]} = \theta_i^2 + (1-\theta_i^2) \frac{\sigma_{i,v}^2}{\sigma_{i,u}^2 + \sigma_{i,v}^2}.
\end{equation}

\subsection{A more general process for productivity growth}

Let $\hat{x}_t \equiv \log x_t - \mu $  denote the demeaned log productivity growth. 
We assume that $\hat{x}_t$ follows the process:
\begin{align}
    \hat{x}_{t+1} &= z_{t} + \sigma_x \left[\sqrt{1 - \rho_{xz}^2} u_{t+1} + \rho_{xz} v_{t+1}\right]\\
    z_{t+1} &= \theta_z z_{t} + \sigma_{z} v_{t+1},
\end{align}
where $u_{t}$ and $v_t$ are standard normal random variables, serially uncorrelated, and uncorrelated with each other. 
Notice that $\mathbb{E}_{t}[x_{t+1}] = z_t$, so $z_t$ corresponds to expected productivity growth. 
The disturbance $v_{t+1}$ can then be interpreted as expectations shocks. 
These expectations shocks are potentially correlated with cash-flow shocks, with correlation coefficient $\rho_{xz}$. 
In the long-run risk literature, $v_{t+1}$ is referred to as a long-run risk shock, while $u_{t+1}$ corresponds to a short-run risk shock.


\paragraph{The ARMA(1,1) case.} Suppose $\rho_{xz} = 1$. This implies that the process for $\hat{x}_t$ specializes to
\begin{equation}
    \hat{x}_{t+1} = \theta_z \hat{x}_t + \sigma_x v_{t+1} - (\theta_z \sigma_x - \sigma_z) v_t,
\end{equation}
which is an ARMA(1,1) process. If we further assume that $\sigma_z = \theta_z \sigma_x$, then we obtain an AR(1) process.

Notice that we can write the conditional expectation of $x_{t+1}$ as follows
\begin{equation}
    \mathbb{E}_t[x_{t+1}] = \theta_z \hat{x}_t - b \frac{\hat{x}_t - \mathbb{E}_{t-1}[x_{t}]}{\sigma_x} \Rightarrow
    \mathbb{E}_t[x_{t+1}] = \frac{\theta_z - b/\sigma_x}{1- b L} \hat{x}_t.
\end{equation}
where $b \equiv \theta_z \sigma_x - \sigma_z$ and $L$ is the lag operator. 

Define $\hat{w}_t$ as follows
\begin{equation}
    \hat{w}_t \equiv \frac{\hat{x}_t}{1-bL} = \sum_{j=1}^\infty b^j \hat{x}_{t-j}.
\end{equation}

\paragraph{Unconditional moments.} The unconditional variance of $z_{t+1}$ is given by 
\begin{equation}
    Var[z_{t+1}] = \frac{\sigma_z^2}{1-\theta_z^2},
\end{equation}
and the unconditional variance of $\hat{x}_{t+1}$ is given by
\begin{equation}
    Var[\hat{x}_{t+1}] = \mathbb{E}\left[Var_t[\hat{x}_{t+1}] \right] + Var[ \mathbb{E}_t[\hat{x}_{t+1}] ] = \sigma_x^2 + \frac{\sigma_z^2}{1-\theta_z^2}. 
\end{equation}

In this general case, we can choose $\theta_z$ and $\sigma_z$ to match the persistence and variance of expectations and choose $\sigma_x^2$ to match the unconditional variance of productivity growth. The parameter $\rho_{xz}$ controls the correlation between expected and realized productivity growth. 

In the special case $\rho_{xz} = 1$ and $\sigma_z = \theta_z \sigma_x$. 
This allows us to match the persistence of expectations and either the unconditional variance of expected productivity growth or unconditional variance of realized productivity growth. 

If $\sigma_z = \theta_z \sigma_x$, then
\begin{equation}
    Var[z_{t+1}] = \theta_z^2\frac{\sigma_x^2}{1-\theta_z^2} = \theta_z^2 Var[x_{t+1}].
\end{equation}

\paragraph{Discretization of the productivity growth process.}  
Using the process for $z_{t+1}$ to eliminate $v_{t+1}$ from the expression for $\hat{x}_t$, we obtain
\begin{equation}
    \hat{x}_{t+1} = \theta_z^{-1} z_{t+1}  + \left(\frac{\rho_{xz} \sigma_x}{\sigma_z} - \frac{1}{\theta_z}\right)\left(  z_{t+1} - \theta_z z_t \right) + \sigma_x  \sqrt{1 - \rho_{xz}^{2}} u_{t+1}.
\end{equation}

Given $z_{t}$, $z_{t+1}$, and $u_{t+1}$, this allow us to solve for $\hat{x}_{t+1}$.
Suppose $z_t$ takes on $N_z$ discrete values and $u_t$ takes on $N_u$ values. 
This implies that $\hat{x}_t$ can take on $N \equiv N_z^2 \times N_u$ values. 
If we impose the constraint $\rho_{xz} = 1$, then $\hat{x}_{t+1}$ is independent of $u_{t+1}$, so $\hat{x}_t$ takes on $N_z^2$ possible values. If we further assume that $\sigma_z = \theta_z \sigma_x$, then $\hat{x}_t$ can take only $N_z$ values. 

The current value of $\hat{x}$ is determined by $(z_{t-1}, z_{t}, u_t) = (z^i, z^j, u^k)$, where $i, j \in \{1, \ldots, N_z\}$ and $k\in \{1, \ldots, N_u\}$. 
We can define the current state as a function of $(i,j,k)$: $s = i + (j-1) N_z + (k-1) N_z^2$. 
The transition matrix is then given by
\begin{equation}
    Pr(s' = i' + (j'-1) N_z + (k'-1) N_z^2 | s) =   \left\{ 
    \begin{matrix}
    Pr(z' = z^{j'}|z = z^j) Pr(u' = u^k), \qquad    &\text{if } i'= j\\
    0, \qquad    &\text{if } i' \neq j
    \end{matrix}
    \right.,
\end{equation}
where $s = i + (j-1) N_z + (k-1) N_z^2$.

We can write the system above in matrix form:
\begin{equation}
    \left[ 
    \begin{matrix}
        \hat{x}_{t+1} \\
        z_{t+1}
    \end{matrix}
    \right] = \left[ 
        \begin{matrix}
            0 & 1 \\
            0 & \theta_z
        \end{matrix}  
    \right]\left[
    \begin{matrix}
        \hat{x}_{t}\\
        z_t
    \end{matrix}
    \right] + \left[ 
    \begin{matrix}
        \sigma_x & 0\\
        \rho_{xz} \sigma_z & \sqrt{1 - \rho_{xz}^2}\sigma_{z}\\
    \end{matrix}
    \right] \left[
    \begin{matrix}
        u_{t+1}\\
        v_{t+1}
    \end{matrix}
    \right].
\end{equation}

Notice that the spectral decomposition of the matrix of coefficients is given by
\begin{equation}
     \left[ 
        \begin{matrix}
            0 & 1 \\
            0 & \theta_z
        \end{matrix}  
    \right]  =
     \left[ 
    \begin{matrix}
        1 & 1 \\
        \theta_z & 0
    \end{matrix}
    \right]
    \left[ 
    \begin{matrix}
        \theta_z & 0\\
        0 & 0
    \end{matrix}
    \right]  \left[ 
    \begin{matrix}
        0 & \theta_z^{-1} \\
        1 & -\theta_z^{-1}
    \end{matrix}
    \right].
\end{equation}

Define the following transformed variables:
\begin{equation}
    \left[
    \begin{matrix}
        w_{1,t}\\
        w_{2,t}
    \end{matrix}
    \right] \equiv \left[ 
    \begin{matrix}
        0 & \theta_z^{-1} \\
        1 & -\theta_z^{-1}
    \end{matrix}
    \right]\left[
    \begin{matrix}
        x_{t}-\mu\\
        z_t
    \end{matrix}
    \right].
\end{equation}

The difference equation for $w_t$ is given by
\begin{equation}
    \left[
    \begin{matrix}
        w_{1,t+1}\\
        w_{2,t+1}
    \end{matrix}
    \right] = \left[ 
    \begin{matrix}
        \theta_z & 0\\
        0 & 0
    \end{matrix}
    \right] 
    \left[
    \begin{matrix}
        w_{1,t}\\
        w_{2,t}
    \end{matrix}
    \right] +
     \left[ 
    \begin{matrix}
        \rho_{xz} \frac{\sigma_{z}}{\theta_z} & \sqrt{1 - \rho_{xz}^2}\frac{\sigma_{z}}{\theta_z}\\
          \sigma_x - \rho_{xz} \frac{\sigma_z }{\theta_z}& -\sqrt{1 - \rho_{xz}^2}\frac{\sigma_{z}}{\theta_z}
    \end{matrix}
    \right]
\left[
    \begin{matrix}
        u_{t+1}\\
        v_{t+1}
    \end{matrix}
    \right].
\end{equation}

We can write the original variables in terms of $w_{1,t}$ and $w_{2,t}$:
\begin{equation}
    \log x_t = \mu + w_{1,t} + w_{2,t}, \qquad \qquad 
    z_t = \theta_z w_{1,t}
\end{equation}

We can then discretize $w_{1,t}$ and $w_{2,t}$.
\begin{equation}
    w_{2,t+1} = \sqrt{1 - \rho_{xz}^2}\sigma_{x} u_{t+1} + (\rho_{xz} \sigma_x - \sigma_z \theta_z^{-1} )v_{t+1}
\end{equation}
\begin{equation}
    w_{1,t+1} = \theta_z w_{1,t} + \sigma_z \theta_z^{-1} u_{t+1}
\end{equation}

\begin{equation}
    x_{t+1} = \mu + w_{1,t+1} + \sqrt{1 - \rho_{xz}^2}\sigma_{x} (w_{1,t+1}  - \theta_z w_{1,t}  ) \frac{\theta_z}{\sigma_z} + (\rho_{xz} \sigma_x - \sigma_z \theta_z^{-1} )v_{t+1}.
\end{equation}




\section{Proofs}

\subsection{Proof of Proposition \ref{prop: oa_Nstate}}\label{proof: oa_Nstate}
\begin{proof}
We provide the characterization of the economy with $N$-possible states in steps, proceeding from the households' problem to the market clearing conditions. 

\paragraph{Step 1: households' problem.} 
Household $i$ chooses consumption $C_{i}$, hours $h_{i}$, and arrow securities  $B_{i}(X,s,s')$ to maximize \eqref{eq: preferences} subject to the budget constraint:
\begin{equation}
    C_{i} + \mathbb{E}_s[ \Lambda(X,s,s') B_{i}(X,s,s')] = B_{i} + W h_{i},
\end{equation}
and an appropriate No-Ponzi condition. 

As in the two-state case, it is useful to transform this budget constraint in terms of net consumption and total wealth:
 \begin{equation}
    \tilde{C}_{i} + \mathbb{E}_s\left[ \Lambda' \left(  B_i' + W' h_i' - \xi' \frac{(h_i')^{1+\nu}}{1+\nu} + \mathcal{H}_i'\right)\right] = B_{i} + W h_i - \xi \frac{h_i^{1+\nu}}{1+\nu} + \mathcal{H}_i \equiv N_i,
\end{equation}
where we used the fact that $\mathcal{H}_i = \mathbb{E}_s\left[ \Lambda' \left( W' h_i' - \xi' \frac{(h_i')^{1+\nu}}{1+\nu} + \mathcal{H}_i'\right)\right]$ and $\tilde{C}_i = C_i - \xi \frac{h_i^{1+\nu}}{1+\nu}$.

We can then write the budget constraint above as follows
\begin{equation}
    \tilde{C}_i + \mathbb{E}_s\left[ \Lambda' N'_i \right] = N_i.
\end{equation}

The household's problem can then be equivalently expressed as choosing $(\tilde{C}_i(N, X,s),N_i'(N,X,s,s') )$ to maximize \eqref{eq: preferences} subject to the constraint above and the natural borrowing limit $N_i'(N,X,s,s') \geq 0$. The solution takes the form in Equation \eqref{eq: value_function}.  It will be useful to define the consumption-wealth ratio $c_i \equiv \frac{\tilde{C}_i}{N}$ and the normalized net worth $n'_i \equiv \frac{N'_i}{N_i}$. Define the portfolio return as $R_{i,n}(X,s,s') \equiv \frac{n'_i(X,s,s')}{1-c_i(X,s)} $, which gives the budget constraint $n' = R_{i,n}'(1-c)$. The function $v_i(X,s)$ must then satisfy the condition
\begin{equation}
    \frac{(v_i(X,s) N)^{1-\psi^{-1}}-1}{1-\psi^{-1}} = \max_{c_{i},n_{i}'}  (1-\beta) \frac{(c_i N)^{1-\psi^{-1}}-1}{1-\psi^{-1}} + \beta \frac{\mathbb{E}_i\left[ (v_i(X',s')n' N)^{1-\gamma}\right]^{\frac{1-\psi^{-1}}{1-\gamma}}-1}{1-\psi^{-1}},
\end{equation}
subject to $n' = R_n'(1-c_i) $, $\mathbb{E}_s[\Lambda' R_n'] = 1$, and $n'\geq 0$. 

\paragraph{Step 2: optimality conditions.} The first-order conditions for the consumption-wealth ratio and the portfolio share are given by 
\begin{align}
    (1-\beta) c_i^{-\psi^{-1}} &= \beta \mathcal{R}_i(X,s)^{1-\psi^{-1}} (1-c_i)^{-\psi^{-1}}\\
    p_{ss'}^i v_i(X',s')^{1-\gamma} R_{i,n}'(X,s,s')^{-\gamma} & = p_{ss'}\Lambda(X,s,s') \mu(X,s),
\end{align}
where $\mathcal{R}_i(X,s) = \mathbb{E}_i\left[ (v_i(X',s') R_{i,n}(X,s,s'))^{1-\gamma} |X,s\right]^{\frac{1}{1-\gamma}}$ and
$\mu(X,s)$ is the (normalized) multiplier on the constraint on returns. From the first-order condition for consumption, we obtain Equation \eqref{eq: c_n}. The envelope condition is given by
\begin{equation}
     v_i(X)^{1-\frac{1}{\psi}} = (1-\beta) c_i^{-\frac{1}{\psi}} \Rightarrow v_i(X)^{1-\gamma} = (1-\beta)^\theta c_i^{-\frac{\theta}{\psi}}.
\end{equation}

Notice that the multiplier is given by
\begin{equation}
    \mu(X,s) = \mathbb{E}_i[ (v_i(X,s') R_{i,n}'(X,s,s'))^{1-\gamma} ] =  \mathcal{R}_{i}(X,s)^{1-\gamma} = \left( \frac{1-\beta}{\beta} \right)^\theta \left[ \frac{c_i}{1-c_i} \right]^{-\frac{\theta}{\psi}}.
\end{equation}

Combining the previous two expressions above with the first-order condition for $R_{i,n}'$, we obtain
\begin{align}
    p_{ss'} \Lambda(X,s,s') &= p_{ss'}^i \frac{(1-\beta)^\theta (c_i')^{-\frac{\theta}{\psi}} R_{i,n}'(X,s,s')^{-\gamma}}{\left( \frac{1-\beta}{\beta} \right)^\theta \left[ \frac{c_i}{1-c_i} \right]^{-\frac{\theta}{\psi}}} \\
    &= p_{ss'}^i \beta^\theta \left( \frac{c_i' N'}{c_i N} \right)^{-\frac{\theta}{\psi}} (R_{i,n}')^{-(1-\theta)}, \\
    &\equiv p_{ss'}^i \Lambda_i(X,s,s'),
\end{align}
using the fact that $\frac{\theta}{\psi} + 1-\theta = \gamma$.

Hence, expressions \eqref{eq: c_n} and \eqref{eq: eulers} hold unchanged with multiple states. Moreover, the change-of-measure equation $\Lambda_i(X,s,s') = \frac{p_{ss'}}{p_{ss'}^i}\Lambda(X,s,s')$ also holds. 

\paragraph{Step 3: firms' problem and labor market outcomes.} 
The firm's problem is essentially the same and the first-order condition \eqref{eq: labor_demand} holds without change. 
The equations for hours and wages \eqref{eq: hours_wages} are also unchanged.

\paragraph{Step 4: law of motion of aggregate state variables.} The aggregate state variables are the same as before. The law of motion of $\mathcal{L}$ is given by
\begin{equation}
    \mathcal{L}'(X,s) = \sum_{s'\in \mathcal{S}} \frac{p_{ss'} \Lambda(X,s,s')}{\sum_{\tilde{s} \in \mathcal{S}} p_{s\tilde{s}} \Lambda(X,s,\tilde{s})} x_{s'},
\end{equation}
and the law of motion of $\eta_i$ is unchanged. 

\paragraph{Step 5: market clearing conditions.}
Notice that $\sum_{i=1}^I \mu_i B_{i}$ must coincide with the cum-dividend value of the firm. Hence, $\sum_{i=1}^I \mu_i N_{i}$ coincides with the cum-dividend value of the surplus claim, $A_{-} P(X,s)$, where $A_{-}$ denotes lagged productivity.

The market clearing condition for net consumption is then given by
\begin{equation}
    \sum_{i=1}^I \mu_i N_i c_i(X,s) = A_{-} \left(x_s h(E)^\alpha - \xi \frac{h(E)^{1+\nu}}{1+\nu}\right).
\end{equation}

Using the fact that $\sum_{i=1}^I \mu_i N_i = A_{-} P(X,s)$, we obtain the market clearing for consumption in Equation \eqref{eq: mkt_clearing_markov}. 
The market clearing for Arrow securities is given by
\begin{equation}
    \sum_{i=1}^I \mu_i N_i n_i(X,s,s') = x_s A_{-} P(X',s').
\end{equation}

We can write the expression above in terms of portfolio returns:
\begin{equation}
    \sum_{i=1}^I \frac{\mu_i N_i(1-c_i(X,s))}{\sum_{j=1}^I \mu_j N_j(1-c_j(X,s))} R_{n,i}(X,s,s') = \frac{x_s A_{-} P(X',s')}{A_{-} \left[ P(X,s) - \left(x_s h(E)^\alpha - \xi \frac{h(E)^{1+\nu}}{1+\nu}\right) \right] },
\end{equation}
using the fact that $\sum_{j=1}^I \mu_j N_j(1-c_j(X,s)) =A_{-} \left[ P(X,s) - \left(x_s h(E)^\alpha - \xi \frac{h(E)^{1+\nu}}{1+\nu}\right) \right] $.

We can write the expression above as follows
\begin{equation}
    \sum_{i=1}^I \tilde{\eta}_i R_{n,i}(X,s,s') = R_p(X,s,s'),
\end{equation}
for each $s' \in \mathcal{S}$.
\end{proof}

\subsection{Proof of Proposition \ref{oa_log_utility}}\label{proof:oa_log_utility}

\begin{proof}
We provide next a characterization of the economy with log-utility and an arbitrary number of states. 

\paragraph{Step 1: consumption and portfolio decisions.} Suppose $\psi = \gamma = 1$. This implies that $c_i(X,s) = 1-\beta$ and that $\Lambda_i(X,s,s') = R_{i,n}^{-1}(X,s,s')$. From the change-of-measure equation, we obtain
\begin{equation}
    R_{i,n}^{-1}(X,s,s') = \frac{p_{ss'} \Lambda(X,s,s')}{p_{ss'}^i} \Rightarrow R_{i,n}(X,s,s') = \frac{p_{ss'}^i}{p_{ss'} \Lambda(X,s,s')}.
\end{equation}

\paragraph{Step 2: the economy's SDF.} 
Plugging the expression for $R_{i,n}(X,s,s')$ into the market clearing condition for Arrow securities paying off in state $s'$, we obtain
\begin{equation}
    \frac{\sum_{i=1}^I \eta_i p_{ss'}^i}{p_{ss'} \Lambda(X,s,s')} = R_p(X,s,s') \Rightarrow
    \Lambda(X,s,s') = \frac{p_{ss'}(X)}{p_{ss'}} R_p^{-1}(X,s,s').
\end{equation}

Notice that the portfolio return for household $i$ is given by
\begin{equation}
    R_{i,n}(X,s,s') = \frac{p_{ss'}^i}{p_{ss'}(X)} R_p(X,s,s').
\end{equation}
Hence, optimistic investors, i.e. investors satisfying $p_{ss'}^i > p_{ss'}(X)$, hold a levered position on the surplus claim. 

\paragraph{Step 3: the surplus claim.} 
From the market clearing condition for goods, we obtain
\begin{equation}
    P(X,s) = \frac{x_s h(E)^\alpha - \xi \frac{h(E)^{1+\nu}}{1+\nu}}{1-\beta}.
\end{equation}

This implies that the return on the surplus claim is given by
\begin{equation}
    R_p(X,s,s') = \frac{x_s P(X',s')}{P(X,s) - \left(x_s h(E)^\alpha - \xi \frac{h(E)^{1+\nu}}{1+\nu}\right)} = \frac{x_s}{\beta} \frac{x_{s'} h(\mathcal{L}'(X,s))^\alpha - \xi \frac{h(\mathcal{L}'(X,s))^{1+\nu}}{1+\nu}}{x_s h(E)^\alpha - \xi \frac{h(E)^{1+\nu}}{1+\nu}}.
\end{equation}

Using the expression for $h(E)$, we can simplify the expression above
\begin{equation}
    R_p(X,s,s') = \frac{x_s}{\beta} \frac{x_{s'} \mathcal{L}'(X,s)^\frac{\alpha}{1+\nu-\alpha} -\frac{ \alpha}{1+\nu} \mathcal{L}'(X,s)^{\frac{1+\nu}{1+\nu-\alpha}}}{x_{s} E^\frac{\alpha}{1+\nu-\alpha} -\frac{ \alpha}{1+\nu} E^{\frac{1+\nu}{1+\nu-\alpha}}},
\end{equation}
which coincides with the expression for the two-type case. 


\paragraph{Step 4: interest rate and risk premium.} Using the fact that $R_b(X,s)$ is the risk-neutral expectation of $R_p(X,s,s')$ and $\mathcal{L}'(X,s)$ is the risk-neutral expectation of $x_{s'}$, we obtain 
\begin{equation}
    R_b(X,s) = \left(1 - \frac{\alpha}{1+\nu} \right) \frac{x_s}{\beta} \frac{  \mathcal{L}'(X,s)^{\frac{1+\nu}{1+\nu-\alpha}}}{x_{s} E^\frac{\alpha}{1+\nu-\alpha} -\frac{ \alpha}{1+\nu} E^{\frac{1+\nu}{1+\nu-\alpha}}},
\end{equation}
which coincides with the expression for the two-type case.  

The expected return on the surplus claim is given by
\begin{equation}
    \mathbb{E}_s[R_p(X,s,s')] = \frac{x_s}{\beta} \frac{\mathbb{E}_s[x_{s'}] \mathcal{L}'(X,s)^\frac{\alpha}{1+\nu-\alpha} -\frac{ \alpha}{1+\nu} \mathcal{L}'(X,s)^{\frac{1+\nu}{1+\nu-\alpha}}}{x_{s} E^\frac{\alpha}{1+\nu-\alpha} -\frac{ \alpha}{1+\nu} E^{\frac{1+\nu}{1+\nu-\alpha}}}.
\end{equation}

Taking the difference of the previous two equations, we obtain the risk premium on the surplus claim:
\begin{equation}
\mathbb{E}_s[R_p^e(X,s,s')] = \frac{x_s}{\beta} \frac{[\mathbb{E}_s[x_{s'}]-\mathcal{L}'(X,s)]\mathcal{L}'(X,s)^\frac{\alpha}{1+\nu-\alpha} }{x_{s} E^\frac{\alpha}{1+\nu-\alpha} -\frac{ \alpha}{1+\nu} E^{\frac{1+\nu}{1+\nu-\alpha}}}.
\end{equation}
\paragraph{Step 5: law of motion of aggregate state variables.}
The risk-neutral probability is given by
\begin{equation}
    \frac{p_{ss'} \Lambda(X,s,s')}{\sum_{\tilde{s}\in \mathcal{S} } p_{s\tilde{s}} \Lambda(X,s,\tilde{s})} = 
    \frac{p_{ss'}(X) R_p^{-1}(X,s,s')}{\sum_{\tilde{s}\in \mathcal{S}} p_{s\tilde{s}}(X) R_p^{-1}(X,s,\tilde{s})} = 
    \frac{p_{ss'}(X) \left[x_{s'} -\frac{ \alpha}{1+\nu} \mathcal{L}'(X,s) \right]^{-1}}{\sum_{\tilde{s}\in \mathcal{S}} p_{s\tilde{s}}(X)\left[x_{\tilde{s}} -\frac{ \alpha}{1+\nu} \mathcal{L}'(X,s) \right]^{-1}}.
\end{equation}

From the law of motion of $\mathcal{L}$, we obtain
\begin{equation}
    \mathcal{L}'(X,s) = \sum_{s'\in \mathcal{S}} x_{s'} \frac{p_{ss'}(X) \left[x_{s'} -\frac{ \alpha}{1+\nu} \mathcal{L}'(X,s) \right]^{-1}}{\sum_{\tilde{s}\in \mathcal{S}} p_{s\tilde{s}}(X)\left[x_{\tilde{s}} -\frac{ \alpha}{1+\nu} \mathcal{L}'(X,s) \right]^{-1}}.
\end{equation}

Rearranging the expression above, we obtain
\begin{equation}
       \sum_{s'\in \mathcal{S}} \frac{p_{ss'}(X)  (x_{s'} - \mathcal{L}'(X,s))}{x_{s'} -\frac{ \alpha}{1+\nu} \mathcal{L}'(X,s) } = 0
\end{equation}

The left-hand side is positive for $\mathcal{L}'(X,s) = x^1$, it is negative for $\mathcal{L}'(X,s) = x^N$, and it is strictly decreasing in $\mathcal{L}'(X,s)$, assuming the condition $x^N < \frac{x^1}{\alpha}$ such that the denominator is positive in the range $x^1 < \mathcal{L}'(X,s) < x^N$. Therefore, a solution exists and it is unique. 

The law of motion of the wealth share is given by
\begin{equation}
    \eta_i'(X,s,s') = \frac{\eta_i R_{i,n}(X,s,s')}{\sum_{j=1}^I\eta_j R_{j,n}(X,s,s')} = \eta_i \frac{R_{i,n}(X,s,s')}{R_{p}(X,s,s')} = \eta_i \frac{p_{ss'}}{p_{ss'}(X)}.
\end{equation}

\end{proof}

\subsection{Proof of Proposition \ref{oa_iid}} \label{proof: oa_iid}
\begin{proof}
We will construct an equilibrium that has iid returns for any financial asset. We guess-and-verify that the consumption-wealth ratio and the net-worth multiplier are constant. 

\paragraph{Step 1: consumption and portfolio decisions.} Let the consumption-wealth ratio be given by $c_i^*(X,s) = 1-\beta^*$, given a constant $\beta^*$ that we need to determine. Given that there is no heterogeneity in beliefs, we obtain from the market clearing condition for Arrow securities $R_{i,n}^*(X,s,s') = R_p^*(X,s,s')$. Plugging $c_i^*(X,s)$ and $R_{i,n}^*(X,s,s')$ into the expression for $\Lambda_i^*(X,s,s')$, we obtain
\begin{equation}
    \Lambda_i^*(X,s,s') = \beta^\theta (\beta^*)^{-\frac{\theta}{\psi}} [R_{p}^*(X,s,s')]^{-\gamma}.
\end{equation}

\paragraph{Step 2: net-worth multiplier.} From the envelope condition, we obtain
\begin{equation}
    v_i^*(X,s)^{1-\psi^{-1}} = (1-\beta) c_i^*(X,s)^{-\psi^{-1}} \Rightarrow v_i^{*}(X,s) = (1-\beta)^{\frac{1}{1-\psi^{-1}}} (1-\beta^*)^{-\frac{\psi^{-1}}{1-\psi^{-1}}}.
\end{equation}

\paragraph{Step 3: wages, hours, and profits.} Using $\alpha = \hat{\alpha} \epsilon$, $\xi = \hat{\xi} \epsilon$, and taking the limit of the expressions for wages, hours, and profits as $\epsilon \rightarrow 0$, we obtain the expressions provided in the proposition.

\paragraph{Step 4: The price and return on the surplus claim.}
For an arbitrary $\alpha$ and $\xi$, the market clearing condition for goods implies that
\begin{equation}
    P^*(X,s) = \left.\frac{x_{s} E^\frac{\alpha}{1+\nu-\alpha} -\frac{ \alpha}{1+\nu} E^{\frac{1+\nu}{1+\nu-\alpha}}}{1-\beta^*}\right|_{\epsilon  = 0} = \frac{x_s}{1-\beta^*}.
\end{equation}

The return on the surplus claim is given by
\begin{equation}
    R_p^*(X,s,s') = \left.\frac{x_{s}}{\beta^*} \frac{x_{s'} \mathcal{L}'(X,s)^\frac{\alpha}{1+\nu-\alpha} -\frac{ \alpha}{1+\nu} \mathcal{L}'(X,s)^{\frac{1+\nu}{1+\nu-\alpha}}}{x_{s} E^\frac{\alpha}{1+\nu-\alpha} -\frac{ \alpha}{1+\nu} E^{\frac{1+\nu}{1+\nu-\alpha}}} \right|_{\epsilon =0} = \frac{x_{s'}}{\beta^*}.
\end{equation}

\paragraph{Step 4: The economy's SDF.}
From the pricing equation, we obtain
\begin{equation}
    \mathbb{E}_i[\Lambda_i^*(X,s,s') R^*_p(X,s,s')] = 1 \Rightarrow \beta^\theta (\beta^*)^{-\theta} \mathbb{E}_i[(x_s')^{1-\gamma}] =1
\end{equation}

Rearranging the expression above, we obtain
\begin{equation}
    \beta^* = \beta \mathbb{E}_i[(x_s')^{1-\gamma}]^{\frac{1-\psi^{-1}}{1-\gamma}}.
\end{equation}

Notice that the condition $\beta^* < 1$ is required to ensure that the consumption-wealth ratio is positive.

The SDF is then given by
\begin{equation}
    \Lambda^*(X,s,s') = \beta \mathbb{E}^*[ x_{s'}^{1-\gamma} ]^{\frac{\gamma - \psi^{-1}}{1-\gamma}} (x_{s'})^{-\gamma},
\end{equation}
using the fact that $\Lambda^*(X,s,s') = \Lambda_i^*(X,s,s')$.

\paragraph{Step 6: Law of motion of aggregate state variables.} The risk-neutral probability is given by
\begin{equation}
    \frac{p_{s'}^* \Lambda^*(X,s,s')}{\sum_{\tilde{s}\in\mathcal{S}} p_{\tilde{s}}^* \Lambda^*(X,s,\tilde{s})} = \frac{p_{s'}^* x_{s'}^{-\gamma}}{\mathbb{E}^*[x_{s'}^{-\gamma}]}
\end{equation}
Hence, $\mathcal{L}'(X,s)$ is given by
\begin{equation}
    \mathcal{L}'(X,s) = \frac{\mathbb{E}^*[x_{s'}^{1-\gamma}]}{\mathbb{E}^*[x_{s'}^{-\gamma}]}.
\end{equation}

\paragraph{Step 5: The interest rate and risk premium on surplus claim.} The interest rate and risk premium are given by
\begin{equation}
    R_b^*(X,s) = \frac{\mathcal{L}'(X,s)}{\beta^*}, \qquad \qquad R_p^e(X,s) = \frac{\mathbb{E}^*[x_{s'}]- \mathcal{L}'(X,s)}{\beta^*}.
\end{equation}

Using the expression for $\beta^*$ and $\mathcal{L}'(X,s)$, we can write the interest rate as follows
\begin{equation}
    R^*_b(X,s) = \beta^{-1} \mathbb{E}^*[ (x_s')^{1-\gamma} ]^{\frac{\psi^{-1}-\gamma}{1-\gamma}} \mathbb{E}^*[(x_{s'})^{-\gamma}]^{-1},
\end{equation}

The expected return on the surplus claim is given by
\begin{equation}
    \mathbb{E}^*\left[R^*_p(X,s,s') \right]= \beta^{-1} \mathbb{E}^*[(x_s')^{1-\gamma}]^{\frac{\psi^{-1}-1}{1-\gamma}} \mathbb{E}^*[x_{s'}]
\end{equation}
and the risk premium on the surplus claim is given by
\begin{equation}
    \mathbb{E}^*\left[\frac{R^*_p(X,s,s')}{R_b^*(X,s)} \right] = \frac{\mathbb{E}^*[x_{s'}] \mathbb{E}^*[x_{s'}^{-\gamma}]}{\mathbb{E}^*[x_{s'}^{1-\gamma}] }
\end{equation}
\end{proof}

\subsection{Proof of Proposition \ref{oa_investor_perturbation}}\label{proof: oa_investor_perturbation}
\begin{proof}
We provide a characterization of the first-order correction for the value function, summarized by the net-worth multiplier $v_i(X,s;\epsilon)$, and the policy functions, namely the consumption-wealth ratio $c_i(X,s;\epsilon)$ and the portfolio return $R_{i,n}(X,s,s';\epsilon)$, given the expansion for the economy's SDF
\begin{equation}
    \Lambda(X,s,s';\epsilon) = \Lambda^*(X,s,s') + \hat{\Lambda}(X,s,s') \epsilon + \mathcal{O}(\epsilon),
\end{equation}
where $\hat{\Lambda}(X,s,s')$ is the first-order correction for the SDF. We take $\hat{\Lambda}(X,s,s')$ as given for now and we will solve for it in a later stage. 

\paragraph{Step 1: value function.} The Bellman equation for household $i$ can be written as follows: 
\begin{align}
    \frac{v_i(X,s;\epsilon)^{1-\psi^{-1}}}{1-\psi^{-1}} &= (1-\beta) \frac{c_i^{1-\psi^{-1}}}{1-\psi^{-1}} + \beta \frac{\left[ \sum_{s'\in \mathcal{S}} p_{ss'}^i(v_i(X',s';\epsilon)R_n' (1-c_i))^{1-\gamma}\right]^{\frac{1-\psi^{-1}}{1-\gamma}}}{1-\psi^{-1}}\\\nonumber
    &+\mu(X,s;\epsilon) \left[ 1 - \sum_{s'\in \mathcal{S}} p_{s'}^* \Lambda(X,s,s';\epsilon) R'_n  \right],
\end{align}
where $X' = \chi(X,s,s';\epsilon)$.  

Taking the derivative of the expression above with respect to $\epsilon$ and evaluating at $\epsilon = 0$, we obtain a condition involving the first-order correction for $v_i$:\small
\begin{align}
  (v_i^*)^{-\psi^{-1}} \hat{v}_i(X,s)  &= \beta \left[ \sum_{s'\in \mathcal{S}} p_{s'}^* 
  (v_i^* R_p^*(s') \beta^*)^{1-\gamma}\right]^{\frac{\gamma-\psi^{-1}}{1-\gamma}} \left[ \sum_{s'\in \mathcal{S}} \delta_{ss'}^i \frac{(v_i^* R_p^*(s') \beta^*)^{1-\gamma}}{1-\gamma} +\right. \nonumber\\
  &\left.\sum_{s'\in \mathcal{S}} p_{s'}^* (v_i^*)^{-\gamma} (R_p^*(s') \beta^* )^{1-\gamma}\hat{v}_i(X^*,s') \right]\
    -\mu^*(X,s) \sum_{s'\in \mathcal{S}} p_{s'}^* \hat{\Lambda}_1(X,s,s') R_p^*(s') ,
\end{align}\normalsize
where we used the fact that $R_n'(X,s,s';0) = R_p^*(X,s,s')$, $c_i(X,s;0) = 1-\beta^*$. We also used the fact that $\chi(X,s,s'; 0) = X^*$, where $X^* = (E^*, \{\eta_i\}_{i=1}^I)$ as  $E'(x,s) = E^*$ and the wealth distribution is constant in the benchmark economy. 

Using the results for the benchmark economy, we can simplify the expression above:
\begin{align}
   \hat{v}_i(X,s)  &= \beta  \left[ \sum_{s'\in \mathcal{S}} p_{s'}^* 
  x_{s'}^{1-\gamma}\right]^{\frac{\gamma-\psi^{-1}}{1-\gamma}} \left[ v_i^* \sum_{s'\in \mathcal{S}} \delta_{ss'}^i \frac{x_{s'}^{1-\gamma}}{1-\gamma} +   \sum_{s'\in \mathcal{S}} p_{s'}^* x_{s'}^{1-\gamma}\hat{v}_i(X^*,s')  \right]\nonumber\\
    &-\mu^*(X,s)(v_i^*)^{\psi^{-1}} \sum_{s'\in \mathcal{S}} p_{s'}^* \hat{\Lambda}_1(X,s,s') R_p(s'),
\end{align}
where we used the fact that $v^*(X,s)$ is constant and $R_p^*(s') \beta^* = x_{s'}$.

Let's solve for $\mu^*(X,s)$ next. The first-order condition for $R_n(X,s,s';\epsilon)$ is given by\small
\begin{equation}
    \beta \left[ \sum_{s'\in \mathcal{S}} p_{ss'}^i(v_i(X',s';\epsilon)R_n' (1-c_i))^{1-\gamma}\right]^{\frac{\gamma-\psi^{-1}}{1-\gamma}} p_{ss'} (v_i(X',s') (1-c_i))^{1-\gamma} R_n(X,s,s')^{-\gamma} = \mu(X,s;\epsilon) p_{s'}^* \Lambda(X,s,s';\epsilon)
\end{equation}\normalsize

Multiplying by $R_n(X,s,s')$ both sides and adding across states, we obtain
\begin{equation}
    \mu(X,s;\epsilon)  = \beta \left[ \sum_{s'\in \mathcal{S}} p_{ss'}^i(v_i(X',s';\epsilon)R_n' (1-c_i))^{1-\gamma}\right]^{\frac{1-\psi^{-1}}{1-\gamma}}.
\end{equation}

Evaluating the expression above at $\epsilon = 0$, we obtain
\begin{equation}
    \mu^*(X,s) = \beta v^*(X,s)^{1-\psi^{-1}} \left[ \sum_{s'\in \mathcal{S}} p_{s'}^*x_{s'}^{1-\gamma}\right]^{\frac{1-\psi^{-1}}{1-\gamma}}.
\end{equation}

Given $\mu^*(X,s)$, we obtain a system of equations for $\hat{v}_i(X,s')$:
\begin{equation}
    \frac{\hat{v}_i(X,s)}{v_i^*(X,s)} - \beta \mathbb{E}^*[x_{s'}^{1-\gamma}]^{\frac{1-\psi^{-1}}{1-\gamma}} \sum_{s'\in \mathcal{S}} \omega^*_{s'} \frac{\hat{v}_i(X^*,s')}{v^*(X,s)} =  \beta \mathbb{E}^*[x_{s'}^{1-\gamma}]^{\frac{1-\psi^{-1}}{1-\gamma}} \left[ \sum_{s'\in \mathcal{S}} \frac{\omega^*_{s'}}{1-\gamma} \frac{\delta_{ss'}^i}{p_{s'}^*} -\sum_{s'\in \mathcal{S}} \omega_{s'}^* \frac{\hat{\Lambda}(X,s,s')}{\Lambda^*(X,s,s')}  \right]
\end{equation}
using the fact that $R_p(s') \Lambda^*(X,s,s') = \frac{x_{s'}^{1-\gamma}}{\mathbb{E}^*[x_{s'}^{1-\gamma}]}$ and the definition $\omega_{s}^* \equiv \frac{p_{s}^* x_s^{1-\gamma}}{\mathbb{E}^*[x_{s'}^{1-\gamma}]}$.

We will solve first for the case $X = X^*$. We can write the system above in matrix form:
\begin{equation}
    \left[
    \begin{matrix}
    1 - \chi_v \omega_{1}^* &
    - \chi_v \omega_{2}^* & \ldots &
    - \chi_v \omega_{N}^* \\
     - \chi_v \omega_{1}^* &
    1- \chi_v \omega_{2}^* & \ldots & 
    - \chi_v \omega_{N}^* \\
    \vdots & \vdots & \ldots & \vdots \\
    - \chi_v \omega_{1}^* &
    - \chi_v \omega_{2}^* & \ldots &
    1- \chi_v \omega_{N}^* 
    \end{matrix}
    \right]
    \left[\begin{matrix}
    \frac{\hat{v}_i(X^*,1)}{v^*(X,s)}\\
    \frac{\hat{v}_i(X^*,2)}{v^*(X,s)}\\
    \vdots\\
    \frac{\hat{v}_i(X^*,N)}{v^*(X,s)}
    \end{matrix}
    \right]
    = \left[\begin{matrix}
    b_{i,1}^v(X^*)\\
    b_{i,2}^v(X^*)\\
    \vdots\\
    b_{i,N}^v(X^*)
    \end{matrix}
    \right],
\end{equation}
where $\chi_v \equiv \beta \mathbb{E}^*[x_{s'}^{1-\gamma}]^{\frac{1-\psi^{-1}}{1-\gamma}}$ and
\begin{equation}
    b_{i,s}^v(X^*) \equiv \chi_v \left[ \sum_{s'\in \mathcal{S}}\frac{ \omega^*_{s'}}{1-\gamma} \frac{\delta_{ss'}^i}{p_{s'}^*} -\sum_{s'\in \mathcal{S}} \omega_{s'}^* \frac{\hat{\Lambda}(X^*,s,s')}{\Lambda^*(X,s,s')}  \right].
\end{equation}

Let $\omega^* = \left[ \omega_1^*,\omega_2^*,\ldots, \omega_N^*\right]$ denote a row vector, $\hat{v}_i(X) = \left[ \hat{v}_i(X, 1), \ldots, \hat{v}_i(X, N)\right]'$ denote a column vector, $b_i^v(X) = \left[ b_{i,1}^v(X), \ldots, b_{i,N}^v(X) \right]' $ denote a column-vector, and $\textbf{1}_N$ denote a $N$-dimensional column vector filled with ones. We can then write the expression above as follows:
\begin{equation}
    \left[ I- \chi_v \mathbb{1}_N\omega^* \right] \frac{\hat{v}_i(X^*)}{v_i^*} = b_{i}^v(X^*).
\end{equation}

The matrix on the left-hand side corresponds to the sum of an invertible matrix and rank-one matrix. An application of the Sherman-Morrison formula gives the inverse of this matrix, which gives the solution
\begin{equation}
    \hat{v}_i(X^*) = v_i^*(X,s)\left[ I + \frac{\chi_v}{1-\chi_v } \mathbb{1}_N \omega^* \right]b_{i}^v(X^*)
\end{equation}

The net-worth multiplier at state $(X,s)$ is then given by
\begin{equation}
    \frac{\hat{v}_i(X^*,s)}{v_i^*(X,s)} = b_{i,s}^v + \frac{\chi_v}{1-\chi_v} \sum_{\tilde{s} \in \mathcal{S}}\omega_{\tilde{s}}^* b_{i,\tilde{s}}^v(X^*).
\end{equation}


Using the expression for $b_{i,s}^v$, we can write the expression above as follows
\begin{equation}
  \frac{\hat{v}_i(X^*,s)}{v_i^*(X,s)} =  \chi_v \sum_{\tilde{s}\in \mathcal{S}}\left(\mathbf{1}_{\tilde{s} =s} + \frac{\chi_v}{1-\chi_v} \omega^*_{\tilde{s}} \right) \left[ \sum_{s'\in \mathcal{S}} \frac{ \omega^*_{s'}}{1-\gamma} \frac{\delta_{\tilde{s}s'}^i}{p_{s'}^*} -\sum_{s'\in \mathcal{S}} \omega_{s'}^* \frac{\hat{\Lambda}(X^*,\tilde{s},s')}{\Lambda^*(X,s,s')}  \right].
\end{equation}

Taking the average of the expression above using the weights $\omega^*_s$, we obtain
\begin{equation}
  \sum_{s\in\mathcal{S}} \omega_s^*\frac{\hat{v}_i(X^*,s)}{v_i^*(X,s)} =  \frac{ \chi_v }{1-\chi_v}\sum_{\tilde{s}\in \mathcal{S}}  \omega^*_{\tilde{s}} \left[ \sum_{s'\in \mathcal{S}} \frac{ \omega^*_{s'}}{1-\gamma} \frac{\delta_{\tilde{s}s'}^i}{p_{s'}^*} -\sum_{s'\in \mathcal{S}} \omega_{s'}^* \frac{\hat{\Lambda}(X^*,\tilde{s},s')}{\Lambda^*(X,s,s')}  \right].
\end{equation}

The net-worth multiplier at $(X,s)$ is then given by
\begin{equation}
    \frac{\hat{v}_i(X,s)}{v_i^*(X,s)}  =  \chi_v \left[ \sum_{s'\in \mathcal{S}} \frac{\omega^*_{s'}}{1-\gamma} \frac{\delta_{ss'}^i}{p_{s'}^*} -\sum_{s'\in \mathcal{S}} \omega_{s'}^* \frac{\hat{\Lambda}(X,s,s')}{\Lambda^*(X,s,s')} +\sum_{s'\in \mathcal{S}} \omega^*_{s'} \frac{\hat{v}_i(X^*,s')}{v^*(X,s)} \right]
\end{equation}

We can then write the expression above as follows:
\begin{equation}
    \frac{\hat{v}_i(X,s)}{v_i^*(X,s)}  =  \chi_v  \sum_{s'\in \mathcal{S}} \omega^*_{s'}\left[\frac{1}{1-\gamma} \frac{\delta_{ss'}^i}{p_{s'}^*} -\frac{\hat{\Lambda}(X,s,s')}{\Lambda^*(X,s,s')} \right] +\chi_v \overline{v},
\end{equation}
where
\begin{equation}
    \overline{v} \equiv \frac{ \chi_v }{1-\chi_v}\sum_{\tilde{s}\in \mathcal{S}}  \omega^*_{\tilde{s}}  \sum_{s'\in \mathcal{S}} \omega^*_{s'}\left[\frac{ 1}{1-\gamma} \frac{\delta_{\tilde{s}s'}^i}{p_{s'}^*} -\frac{\hat{\Lambda}(X^*,\tilde{s},s')}{\Lambda^*(X^*,\tilde{s},s')}  \right].
\end{equation}

\paragraph{Step 2: consumption-wealth ratio.} From the envelope condition, the consumption-wealth ratio is given by
\begin{equation}
    c_i(X,s;\epsilon) = (1-\beta)^\psi v_i(X,s;\epsilon)^{1-\psi}.
\end{equation}

The first-order correction for consumption is then given by
\begin{equation}
    \hat{c}_1(X,s) = (1-\beta)^{\psi} (v_i^*(X,s))^{1-\psi} (1-\psi) \frac{\hat{v}_i(X,s)}{v_i^*(X,s)}.
\end{equation}

\paragraph{Step 3: portfolio return.} 
Using the expression for the Lagrange multiplier, we can write the first-order condition for the portfolio return as follows
\begin{equation}
    \frac{p_{ss'}^i}{p_{s'}^*} v_i(X',s')^{1-\gamma} R_n(X,s,s';\epsilon)^{-\gamma} = \Lambda(X,s,s';\epsilon) \sum_{s'\in \mathcal{S}} p_{ss'}^i(v_i(X',s';\epsilon)R_n(X,s,s';\epsilon) )^{1-\gamma}
\end{equation}

Expanding the expression above in $\epsilon$, we obtain
\begin{align}
    \frac{\delta_{ss'}^i}{p_{s'}^*} &+ (1-\gamma)  \frac{\hat{v}_i(X^*,s')}{v^*(X,s)} -\gamma \frac{\hat{R}_{n,i}(X,s,s')}{R^*_p(X,s,s')} = \frac{\hat{\Lambda}(X,s,s')}{\Lambda^*(X,s,s')}\\
    &+ \sum_{s'\in\mathcal{S}'} \frac{p_{s'}^*(v^*(X,s)R_p^*(s'))^{1-\gamma}}{\sum_{\tilde{s} \in \mathcal{S}} p_{\tilde{s}}^*(v^*(X,s)R_p^*(\tilde{s}))^{1-\gamma}} \left[ \frac{\delta_{ss'}^i}{p_{s'}^*} + (1-\gamma) \left(\frac{\hat{v}_i(X^*,s')}{v^*(X,s)} + \frac{\hat{R}_{n,i}(X,s,s')}{R_p^*(X,s,s')}\right)\right]
\end{align}

Rearranging the expression above, we obtain
\begin{equation}
    \gamma \frac{\hat{R}_{n,i}(X,s,s')}{R^*_p(X,s,s')} +  (1-\gamma) \sum_{\tilde{s}\in\mathcal{S}'} \omega^*_{\tilde{s}} \frac{\hat{R}_{n,i}(X,s,\tilde{s})}{R^*_p(X,s,\tilde{s})} = b_i^R(X,s,s'),
\end{equation}
where
\begin{align}
    b_i^R(X,s,s') &\equiv \frac{\delta_{ss'}^i}{p_{s'}^*} + (1-\gamma)  \frac{\hat{v}_i(X^*,s')}{v^*(X,s)}- \sum_{\tilde{s}\in\mathcal{S}} \omega^*_{\tilde{s}} \left[ \frac{\delta_{s\tilde{s}}^i}{p_{\tilde{s}}^*} + (1-\gamma) \frac{\hat{v}_i(X^*,\tilde{s})}{v^*(X,s)}\right]-\frac{\hat{\Lambda}_1(X,s,s')}{\Lambda^*(X,s,s')} 
\end{align}

We can write the system above in matrix form:
\begin{equation}
    \left[
    \begin{matrix}
    \gamma + (1-\gamma) \omega^*_{1} & (1-\gamma) \omega^*_{2} & \ldots & (1-\gamma) \omega^*{N} \\
    (1-\gamma) \omega^*_{1} & \gamma + (1-\gamma) \omega^*_{2} & \ldots & (1-\gamma) \omega^*_{N} \\
    \vdots & \vdots & \ldots  & \vdots \\
    (1-\gamma) \omega^*_{1} &  (1-\gamma) \omega^*_2 & \ldots & \gamma + (1-\omega) \omega^*_{N}
    \end{matrix}
    \right] \left[ 
    \begin{matrix}
    \frac{\hat{R}_{n,i}(X,s,1)}{R^*_p(X,s,1)} \\
    \frac{\hat{R}_{n,i}(X,s,2)}{R^*_p(X,s,2)} \\
    \vdots\\
    \frac{\hat{R}_{n,i}(X,s,N)}{R^*_p(X,s,N)} 
    \end{matrix}
    \right] = 
    \left[ 
    \begin{matrix}
    b^R_i(X,s, 1) \\
    b^R_i(X,s, 2) \\
    \vdots\\
    b^R_i(X,s, N)
    \end{matrix}
    \right]
\end{equation}

Denote the matrix above by $A^*$ and define the row vector $\omega^* \equiv [\omega^*(s_1), \omega^*(s_2), \ldots, \omega^*(s_N)]$ and the column-vector $\mathbb{1}$ with $1$ in every entry. We can then write $A^*$ as follows:
\begin{equation}
    A^* = \gamma I + (1-\gamma) \mathbb{1} \omega^*. 
\end{equation}

The inverse of $A^*$ is given by
\begin{equation}
    (A^*)^{-1} = \frac{1}{\gamma} I - \frac{1-\gamma}{\gamma} \mathbb{1} \omega^*. 
\end{equation}

The portfolio return is then given by
\begin{equation}
    \frac{\hat{R}_{n,i}(X,s,s')}{R^*_p(X,s,s')} = \frac{1}{\gamma} b_i^R(X,s,s') + \frac{1-\gamma}{\gamma} \sum_{\tilde{s} \in \mathcal{S}} \omega^*_{\tilde{s}} \frac{\hat{\Lambda}_1(X,s,\tilde{s})}{\Lambda^*(X,s,\tilde{s})}.
\end{equation}

We can write the expression above as follows:\small
\begin{equation}
    \frac{\hat{R}_{n,i}(X,s,s')}{R^*_p(X,s,s')} = \frac{1}{\gamma}  \sum_{\tilde{s}\in\mathcal{S}} \omega^*_{\tilde{s}} \left[ \left(\frac{\delta_{ss'}^i}{p_{s'}^*}-\frac{\delta_{s\tilde{s}}^i}{p_{\tilde{s}}^*}\right) -\frac{\hat{\Lambda}_1(X,s,s')}{\Lambda^*(X,s,s')}  \right]  + \frac{1-\gamma}{\gamma} \sum_{\tilde{s} \in \mathcal{S}}\omega^*_{\tilde{s}} \left[
    \left(\frac{\hat{v}_i(X^*,s')}{v^*(X,s)}-\frac{\hat{v}_i(X^*,\tilde{s})}{v^*(X,s)}\right)+
     \frac{\hat{\Lambda}_1(X,s,\tilde{s})}{\Lambda^*(X,s,\tilde{s})}\right].
\end{equation}\normalsize

Notice that we can write the term involving  $\hat{v}_i(X,s)$ as follows\footnotesize
\begin{align}
    \frac{\hat{v}_i(X^*,s')}{v^*(X^*,s')}-\sum_{\tilde{s} \in \mathcal{S}}\omega^*_{\tilde{s}} 
    \frac{\hat{v}_i(X^*,\tilde{s})}{v^*(X^*,\tilde{s})} &= 
    \chi_v \sum_{\tilde{s}\in\mathcal{S}} \omega^*_{\tilde{s}} \sum_{\tilde{s}'\in \mathcal{S}} \omega^*_{\tilde{s}'}\left[\frac{1}{1-\gamma} \left(\frac{\delta_{s'\tilde{s}'}^i}{p_{\tilde{s}'}^*} -\frac{\delta_{\tilde{s}\tilde{s}'}^i}{p_{\tilde{s}'}^*} \right) -\left( 
    \frac{\hat{\Lambda}(X^*,s',\tilde{s}')}{\Lambda^*(X^*,s',\tilde{s}')}-
    \frac{\hat{\Lambda}(X^*,\tilde{s},\tilde{s}')}{\Lambda^*(X^*,\tilde{s},\tilde{s}')}
    \right) \right]
\end{align}\normalsize

\end{proof}

\subsection{Proof of Proposition \ref{oa_labor_mkt}}\label{proof: oa_labor_mkt}

\begin{proof}
From the expression for wages, we obtain:
\begin{equation}
    w(E; \epsilon) = \xi \left( \frac{\hat{\alpha} E}{\hat{\xi}} \right)^{\frac{\nu}{1+\nu -\alpha}} = \hat{\xi}\left( \frac{\hat{\alpha} E}{\hat{\xi}} \right)^{\frac{\nu}{1+\nu}} \epsilon + \mathcal{O}(\epsilon^2).
\end{equation}

Hours are given by
\begin{equation}
    h(E; \epsilon) = \exp \left[\frac{1}{1+\nu -\alpha} \log \left( \frac{\hat{\alpha} E}{\hat{\xi}} \right) \right] = \left( \frac{\hat{\alpha} E}{\hat{\xi}} \right)^{\frac{1}{1+\nu}} + \left( \frac{\hat{\alpha} E}{\hat{\xi}} \right)^{\frac{1}{1+\nu}} \frac{\log \left( \frac{\hat{\alpha} E}{\hat{\xi}} \right)}{(1+\nu)^2} \hat{\alpha}\epsilon + \mathcal{O}(\epsilon^2).
\end{equation}

Profits are given by 
\begin{equation}
    \pi(X,s;\epsilon) = \left( \frac{\hat{\alpha}}{\hat{\xi}} \right)^{\frac{\alpha}{1+\nu -\alpha}} \left[  x_s E^{\frac{\alpha}{1+\nu - \alpha}} - \alpha E^{\frac{1+\nu}{1+\nu - \alpha}}\right] = x_s +  \left[ x_s \frac{\log \left( \hat{\alpha} E/   \hat{\xi} \right)}{1+\nu} -E \right] \hat{\alpha} \epsilon + \mathcal{O}(\epsilon^2),
\end{equation}
where we used the following Taylor expansion:
\begin{align}
    E^{\frac{\alpha}{1+\nu - \alpha}} &= 1 + \frac{\log E}{1+\nu}  \hat{\alpha} \epsilon + \mathcal{O}(\epsilon^2).
\end{align}
\end{proof}

\subsection{Proof of Proposition \ref{oa_asset_prices}}\label{proof: oa_asset_prices}

\begin{proof} We derive next the expression for the price and return for the surplus claim and riskless asset.
\paragraph{Step 1: price of surplus claim.} The market clearing for consumption can be written as
\begin{equation}
    P(X,s;\epsilon) \sum_{i=1}^I \eta_i c_i(X,s;\epsilon) = \left( \frac{\alpha }{\xi} \right)^{\frac{\alpha}{1+\nu-\alpha}} \left[x_s E^{\frac{\alpha}{1+\nu-\alpha}} - \frac{\alpha}{1+\nu}  E^{\frac{1+\nu}{1+\nu-\alpha}}\right].
\end{equation}

Expanding the expression above in $\epsilon$, we obtain
\begin{equation}
    \frac{\hat{P}(X,s)}{P^*(X,s)} + \sum_{i=1}^I \eta_i \frac{\hat{c}_i(X,s)}{c^*(X,s)} = \left[\frac{\log (\hat{\alpha} E / \hat{\xi})}{1+\nu} - \frac{1}{1+\nu} \frac{E}{x_s}\right] \hat{\alpha}.
\end{equation}

Rearranging the expression above, and using the expression for $\hat{c}_i(X,s)$, we obtain
\begin{equation}
    \frac{\hat{P}(X,s)}{P^*(X,s)}  = \left[\log (\hat{\alpha} E / \hat{\xi}) -  \frac{E}{x_s}\right] \frac{\hat{\alpha}}{1+\nu} - (1-\psi) \sum_{i=1}^I \eta_i \frac{\hat{v}_i(X,s)}{v^*(X,s)}.
\end{equation}


\paragraph{Step 2: return on surplus claim.} The return on the surplus claim is defined as follows
\begin{equation}
    R_p(X,s,s';\epsilon) = \frac{x_s P(\chi(X,s,s';\epsilon), s';\epsilon)}{P(X,s;\epsilon) - \left( x_s h(E)^\alpha - \xi \frac{h(E;\epsilon)^{1+\nu}}{1+\nu} \right)}.
\end{equation}

Expanding the expression above in $\epsilon$, we obtain
\begin{equation}
    \frac{\hat{R}_{p}(X,s,s')}{R_p^*(X,s,s')} = \frac{\hat{P}(X^*,s')}{P^*(X^*,s')} - \left[\frac{P^*(X,s)}{P^*(X,s) - x_s } \frac{\hat{P}(X,s)}{P^*(X,s)} -\frac{1}{P^*(X,s) - x_s } \left(  x_s\log\left( \frac{\hat{\alpha}}{\hat{\xi}} E\right)  -E \right)   \frac{\hat{\alpha}}{1+\nu} \right].
\end{equation}

We can write the expression above as follows:
\begin{equation}
    \frac{\hat{R}_{p}(X,s,s')}{R_p^*(X,s,s')} = \frac{\hat{P}(X^*,s')}{P^*(X^*,s')} - \left[(\beta^{*})^{-1} \frac{\hat{P}(X,s)}{P^*(X,s)} +(1-(\beta^{*})^{-1}) \left(\log\left( \frac{\hat{\alpha}}{\hat{\xi}} E\right)  -\frac{E}{x_s} \right)   \frac{\hat{\alpha}}{1+\nu} \right],
\end{equation}
using the fact that $P^*(X,s) = x_s/(1-\beta^*)$.

Using the expression for the price of the surplus claim, we obtain
\begin{equation}
    \frac{\hat{R}_{p}(X,s,s')}{R_p^*(X,s,s')} = 
    \left[ \log(E^*/E) - \left( \frac{E^*}{x_{s'}} - \frac{E}{x_s} \right) \right]\frac{\hat{\alpha}}{1+\nu} - (1-\psi) \sum_{i=1}^I \eta_i \left[ \frac{\hat{v}_i(X^*,s')}{v^*(X^*,s')}- \frac{1}{\beta^*} \frac{\hat{v}_i(X,s)}{v^*(X,s)}\right],
\end{equation}

\paragraph{Step 3: interest rate.} The interest rate is given by
\begin{equation}
    R_b(X,s,s';\epsilon) = \left[ \sum_{s'\in \mathcal{S}} p_{s'}^* \Lambda(X,s,s';\epsilon) \right]^{-1} \Rightarrow
    \frac{\hat{R}_b(X,s)}{R^*_b(X,s)} = - \sum_{s'\in\mathcal{S}} \frac{p_{s'} x_{s'}^{-\gamma}}{\mathbb{E}^*[x_{s'}^{-\gamma}]} \frac{\hat{\Lambda}(X,s,s')}{\Lambda^*(X,s,s')}.
\end{equation}

\paragraph{Step 4: stock prices.} Stock prices, normalized by current productivity, satisfy the functional equation:
\begin{equation}
    Q(X,s;\epsilon) = \sum_{s'\in\mathcal{S}} p_{s'}^* \Lambda(X,s,s'; \epsilon) \left[\left( \frac{\hat{\alpha}}{\hat{\xi}} \right)^{\frac{\alpha}{1+\nu-\alpha}} \left( x_{s'} \mathcal{L}'(X,s;\epsilon)^{\frac{\alpha}{1+\nu - \alpha}} - \alpha \mathcal{L}'(X,s;\epsilon)^{\frac{1+\nu}{1+\nu - \alpha}} \right) + x_{s'} Q(\chi(X,s,s';\epsilon), s';\epsilon) \right]
\end{equation}

For $\epsilon = 0$, we obtain
\begin{equation}
    Q^*(X,s) = \sum_{s'\in\mathcal{S}} p_{s'}^* \Lambda^*(X,s,s') x_{s'}\left[ 1 +  Q^*(X^*, s') \right].
\end{equation}

We can write the expression above as follows:
\begin{equation}
    Q^*(X,s) = \beta^* \sum_{s'\in\mathcal{S}} \frac{p_{s'}^* x_{s'}^{1-\gamma}}{\mathbb{E}^*[x_{s'}^{1-\gamma}]}\left[ 1 +  Q^*(X^*, s') \right] \Rightarrow
    Q^*(X,s) = \frac{\beta^*}{1-\beta^*}.
\end{equation}

Expanding the expression for $Q(X,s)$, we obtain
\begin{equation}
    \frac{\hat{Q}(X,s)}{Q^*(X,s)} = \sum_{s'\in\mathcal{S}} \omega_{s'}^* \left[\frac{\hat{\Lambda}(X,s,s')}{\Lambda^*(X,s,s')} + (1-\beta^*) \left( \frac{\log \left( \frac{\hat{\alpha} E^*}{\hat{\xi}} \right)}{1+\nu} - \frac{E^*}{x_{s'}}\right) \hat{\alpha}  + \beta^* \frac{\hat{Q}(X^*,s')}{Q(X^*,s)} \right].
\end{equation}

Evaluating the expression above at $X = X^*$, we obtain
\begin{equation}
    \left[I - \beta^* \mathbb{1}_N \omega^* \right] \hat{Q}(X^*) = b^Q(X^*),
\end{equation}
where $\hat{Q}(X) \equiv [\hat{Q}(X,1), \ldots, \hat{Q}(X,N)]' 
$, $b^Q(X) \equiv \left[b^Q(X,1),\ldots, b^Q(X,N) \right]'$, and $b^Q(X,s) \equiv   Q^*(X,s) \sum_{s'\in\mathcal{S}} \omega_{s'}^* \left[ \frac{\hat{\Lambda}(X,s,s')}{\Lambda^*(X,s,s')} + (1-\beta^*) \left( \frac{\log \left( \frac{\hat{\alpha} E^*}{\hat{\xi}} \right)}{1+\nu} - \frac{E^*}{x_{s'}}\right) \hat{\alpha} \right]$. 

Solving the system above, we obtain
\begin{equation}
     \hat{Q}(X^*) = \left[ I + \frac{\beta^*}{1-\beta^*} \mathbb{1}_N \omega^*\right]b^Q(X^*).
\end{equation}

We can then write $\hat{Q}(X,s)$ as follows:\footnotesize
\begin{equation}
    \frac{\hat{Q}(X,s)}{Q^*(X,s)} = \sum_{s'\in\mathcal{S}} \omega_{s'}^* \left[\frac{\hat{\Lambda}(X,s,s')}{\Lambda^*(X,s,s')} + \beta^* \frac{\hat{\Lambda}(X^*,s,s')}{\Lambda^*(X^*,s,s')}+  \frac{(\beta^*)^2}{1-\beta^*}  \sum_{\tilde{s}\in\mathcal{S}}\omega^*_{\tilde{s}}\frac{\hat{\Lambda}(X^*,\tilde{s},s')}{\Lambda^*(X^*,s,s')} \right]+ \left( \frac{\log \left( \frac{\hat{\alpha} E^*}{\hat{\xi}} \right)}{1+\nu} - \sum_{s'\in \mathcal{S}}\omega_{s'}^*\frac{E^*}{x_{s'}}\right) \hat{\alpha}.
\end{equation}\normalsize

\paragraph{Step 5: equity returns.} Equity returns are given by 
\begin{equation}
    R_E(X,s,s';\epsilon) = \frac{x_{s'} Q(\chi(X,s,s';\epsilon), s') + \pi(\mathcal{L}'(X,s;\epsilon; \epsilon), s') }{Q(X,s;\epsilon)}.
\end{equation}

Evaluating the expression above at $\epsilon = 0$, we obtain
\begin{equation}
    R_E^*(X,s,s') = \frac{x_{s'} Q^*(X^*, s') + \pi^*(E^*, s') }{Q^*(X,s;\epsilon)} = \frac{x_{s'}}{\beta^*}.
\end{equation}

The first-order correction is given by 
\begin{equation}
    \frac{\hat{R}_E(X,s,s')}{R^*_E(X,s,s')} = \beta^* \frac{\hat{Q}(X^*,s')}{Q^*(X,s)} +(1-\beta^*) \frac{\hat{\pi}(E^*, s')}{x_{s'}} - \frac{\hat{Q}(X,s)}{Q^*(X,s)}.
\end{equation}

\paragraph{Step 6: conditional risk premium.} The conditional risk premium is defined as follows:
\begin{equation}
    \overline{R}_E(X,s;\epsilon) = \sum_{s'\in \mathcal{S}} p_{s'}^*\left[ \frac{R_E(X,s,s';\epsilon)}{R_b(X,s;\epsilon)}\right].
\end{equation}

The first-order correction is given by
\begin{equation}
    \frac{\widehat{\overline{R}}_E(X,s)}{\overline{R}^*(X,s)} = \sum_{s'\in \mathcal{S}} \frac{p_{s'}^* x_{s'}}{\mathbb{E}^*[x_{s'}]} \frac{\hat{R}_E(X,s,s')}{R_E^*(X,s,s')} -\frac{\hat{R}_b(X,s)}{R^*_b(X,s)}.
\end{equation}


\end{proof}

\subsection{Proof of Proposition \ref{oa_aggregate_state}}\label{proof: oa_aggregate_state}
\begin{proof}
We consider next the law of motion of $\eta_i$ and $\mathcal{L}$.
\paragraph{Step 1: wealth distribution.} The law of motion of $\eta_i$ can be written as
\begin{equation}
    \eta_i'(X,s,s';\epsilon) \sum_{j=1}^I\eta_j R_{j,n}(X,s,s';\epsilon) (1-c_j(X,s;\epsilon))  = \eta_i R_{i,n}(X,s,s';\epsilon) (1-c_i(X,s;\epsilon)). 
\end{equation}

Expanding the expression above in $\epsilon$, we obtain
\begin{equation}
    \hat{\eta}_i'(X,s,s') = \eta_i \left[ \frac{\hat{R}_{i,n}(X,s,s')}{R^*_{i,n}(X,s,s')}-\sum_{j=1}^I \eta_i \frac{\hat{R}_{j,n}(X,s,s')}{R^*_{j,n}(X,s,s')} - \frac{c_i^*(X,s)}{1-c_i^*(X,s)}\left(\frac{\hat{c}_i(X,s)}{c^*_i(X,s)}  - \sum_{j=1}^I \eta_j \frac{\hat{c}_j(X,s)}{c_j^*(X,s)} \right)\right].
\end{equation}
Using $c_i^*(X,s) = 1-\beta^*$ gives the expression in the proposition.

\paragraph{Step 2: risk-neutral probability of productivity growth.} The law of motion of $\mathcal{L}$ can be written as
\begin{equation}
    \mathcal{L}'(X,s;\epsilon)  = R_b(X,s;\epsilon) \sum_{s'\in\mathcal{S}}p_{s'}^* \Lambda(X,s,s';\epsilon) x_{s'}.
\end{equation}

Expanding the expression above in $\epsilon$, we obtain
\begin{equation}
    \frac{\hat{\mathcal{L}}'(X,s)}{E^*(X,s)} = \frac{ \hat{R}_b(X,s) }{R_b^*(X,s)} +  \sum_{s'\in\mathcal{S}}\frac{p_{s'}^* \Lambda^*(X,s,s')x_{s'}}{\mathbb{E}^*[\Lambda^*(X,s,s') x_{s'}]} \frac{\hat{\Lambda}(X,s,s')}{\Lambda^*(X,s,s')} 
\end{equation}

We can write the expression above as follows:
\begin{equation}
    \frac{\hat{\mathcal{L}}'(X,s)}{E^*} = \frac{\hat{R}_b(X,s)}{R_b^*(X,s)} +  
    \sum_{s'\in \mathcal{S}} \frac{p_{s'}^* x_{s'}^{1-\gamma}}{\mathbb{E}^*[x_{s'}^{1-\gamma}]}\frac{\hat{\Lambda}(X,s,s')}{\Lambda^*(X,s,s')}.
\end{equation}
Using the definition of $E^*$ and $\omega_{s}^*$, we obtain the expression given in the proposition.
\end{proof}

\subsection{Proof of Proposition \ref{oa_SDF}}\label{proof: oa_SDF}
\begin{proof}
We consider the derivation of the economy's SDF $\hat{\Lambda}(X,s,s')$.
\paragraph{Step 1: the system of equations.}
The market clearing for the Arrow security paying off in state $s'$ is given by
\begin{equation}
    \sum_{i=1}^I \eta_i (1-c_i(X,s;\epsilon)) R_{n,i}(X,s,s') =R_p(X,s,s') \sum_{i=1}^I \eta_i (1-c_i(X,s;\epsilon)).
\end{equation}

Expanding the expression above, we obtain
\begin{equation}
    \sum_{i=1}^I \eta_i \hat{R}_{n,i}(X,s,s') = \hat{R}_{p}(X,s,s').
\end{equation}

Using the expression for $\hat{R}_{n,i}(X,s,s')$ and $\hat{R}_p(X,s,s')$, we obtain\footnotesize
\begin{align}
    \frac{1}{\gamma} &\left[\frac{\delta_{ss'}(X)}{p_{s'}^*} - \sum_{\tilde{s}'\in\mathcal{S}} \omega^*_{\tilde{s}'} \frac{\delta_{s\tilde{s}'}(X)}{p_{\tilde{s}'}^*} - \frac{\hat{\Lambda}(X,s,s')}{\Lambda^*(X,s,s')}  \right]+\frac{1-\gamma}{\gamma} \left[ \sum_{i=1}^I\eta_i\left[\frac{\hat{v}_i(X^*,s')}{v^*(X,s)}-\sum_{\tilde{s} \in \mathcal{S}}\omega^*_{\tilde{s}} 
    \frac{\hat{v}_i(X^*,\tilde{s})}{v^*(X,s)}\right]+
     \sum_{\tilde{s} \in \mathcal{S}}\omega^*_{\tilde{s}} \frac{\hat{\Lambda}(X,s,\tilde{s})}{\Lambda^*(X,s,\tilde{s})}\right] = \nonumber\\
     &\left[ \log\frac{E^*}{E} - \left( \frac{E^*}{x_{s'}} - \frac{E}{x_s} \right) \right]\frac{\hat{\alpha}}{1+\nu} + (\psi-1) \sum_{i=1}^I \eta_i \left[ \frac{\hat{v}_i(X^*,s')}{v^*(X^*,s')}- \frac{1}{\beta^*} \frac{\hat{v}_i(X,s)}{v^*(X,s)}\right].
\end{align}\normalsize

Using the expression for $\hat{v}_i(X,s)$, we obtain\footnotesize
\begin{align}
    &\frac{1}{\gamma}\frac{\hat{\Lambda}(X,s,s')}{\Lambda^*(X,s,s')} + \frac{1-\gamma}{\gamma} \left[ \beta^* \sum_{\tilde{s}\in \mathcal{S}}\omega^*_{\tilde{s}} \sum_{\tilde{s}'\in \mathcal{S}} \omega_{\tilde{s}'}^* \left( \frac{\hat{\Lambda}(X^*,s' ,\tilde{s}')}{\Lambda^*(X^*,s' ,\tilde{s}')} -\frac{\hat{\Lambda}(X^*,\tilde{s},\tilde{s}')}{\Lambda^*(X^*,\tilde{s},\tilde{s}')}   \right)-\sum_{\tilde{s}'  \in \mathcal{S}}\omega^*_{\tilde{s}' } \frac{\hat{\Lambda}(X,s,\tilde{s}' )}{\Lambda^*(X,s,\tilde{s}')} \right]\\
    &(\psi-1) \beta^ *\left[- \sum_{\tilde{s}'\in\mathcal{S}}\omega^*_{\tilde{s}'} \left(\frac{\hat{\Lambda}(X^*,s',\tilde{s}')}{\Lambda^*(X^*,s',\tilde{s}')} - \frac{1}{\beta^*}\frac{\hat{\Lambda}(X,s,\tilde{s}')}{\Lambda^*(X,s,\tilde{s}')}\right) + \sum_{\tilde{s}\in \mathcal{S}}\omega^*_{\tilde{s}} \sum_{\tilde{s}'\in \mathcal{S}} \omega_{\tilde{s}'}^* \frac{\hat{\Lambda}(X^*,\tilde{s},\tilde{s}') }{\Lambda^*(X^*,\tilde{s},\tilde{s}')} \right] =b^\Lambda(X,s,s'),
\end{align}\normalsize
where
\begin{align}
    b^\Lambda&(X,s,s') \equiv \frac{1}{\gamma} \left[\frac{\delta_{ss'}(X)}{p_{s'}^*} - \sum_{\tilde{s}'\in\mathcal{S}} \omega^*_{\tilde{s}'} \frac{\delta_{s\tilde{s}'}(X)}{p_{\tilde{s}'}^*}\right]+\frac{\beta^*}{\gamma} \left[ \sum_{\tilde{s}\in \mathcal{S}}\omega^*_{\tilde{s}} \sum_{\tilde{s}'\in \mathcal{S}} \omega_{\tilde{s}'}^*\left[ \frac{\delta_{s' \tilde{s}'}(X)}{p_{\tilde{s}'}^*}- \frac{\delta_{\tilde{s}\tilde{s}'}(X)}{p_{\tilde{s}'}^*}\right]\right] +\nonumber\\
    &-\frac{\psi-1}{1-\gamma} \sum_{\tilde{s}\in \mathcal{S}}\omega^*_{\tilde{s}} \sum_{\tilde{s}'\in \mathcal{S}} \omega_{\tilde{s}'}^* \left[ \frac{\delta_{s\tilde{s}'}(X)}{p^*_{\tilde{s}'}}-\beta^*\frac{\delta_{s'\tilde{s}'}(X^*)}{p^*_{\tilde{s}'}} + \beta^* \frac{\delta_{\tilde{s}\tilde{s}'}(X)}{p^*_{\tilde{s}'}} \right]  - \left[ \log\frac{E^*}{E} - \left( \frac{E^*}{x_{s'}} - \frac{E}{x_s} \right) \right]\frac{\hat{\alpha}}{1+\nu}.
\end{align}

We can simplify the expression above as follows:\footnotesize
\begin{align}
    &\frac{1}{\gamma}\frac{\hat{\Lambda}(X,s,s')}{\Lambda^*(X,s,s')} +(\psi - \gamma^{-1}) \left[\omega^* \cdot \hat{\Lambda}(X,s) -\beta^* \omega^* \cdot \hat{\Lambda}(X^*,s') + \beta^* \sum_{\tilde{s}\in \mathcal{S}} \omega_{\tilde{s}} ( \omega^* \cdot \hat{\Lambda}(X^*,\tilde{s}) ) \right] = b^\Lambda(X,s,s'),
\end{align}\normalsize
where $\hat{\Lambda}(X,s) = \left[\frac{\hat{\Lambda}(X^*,s,1)}{\Lambda^*(X^*,s,1)}, \frac{\hat{\Lambda}(X^*,s,2)}{\Lambda^*(X^*,s,2)}, \ldots, \frac{\hat{\Lambda}(X^*,s,N)}{\Lambda^*(X^*,s,N)} \right]'$ and $\omega^* \cdot \hat{\Lambda}(X, s) = \sum_{\tilde{s}' } \omega^*_{\tilde{s}'} \frac{\hat{\Lambda}(X,s,\tilde{s}')}{\Lambda^*(X,s,\tilde{s}')}$.

\paragraph{Step 2: solving the system.} We can write the system above in matrix form as follows:
\begin{equation}
    \left[ \gamma^{-1} I + (\psi - \gamma^{-1}) \mathbb{1}_N \omega^*\right] \hat{\Lambda}(X,s) = \tilde{b}^{\Lambda}(X,s),
\end{equation}
where $\tilde{b}^{\Lambda}(X,s) = [\tilde{b}^{\Lambda}(X,s,1),\tilde{b}^{\Lambda}(X,s,2),\ldots, \tilde{b}^{\Lambda}(X,s,N) ]'$ and 
\begin{equation}
    \tilde{b}^\Lambda(X,s,s') = b^{\Lambda}(X,s,s') + (\psi - \gamma^{-1}) \beta^* \left[ \omega^* \cdot \hat{\Lambda}(X^*,s') - \sum_{\tilde{s}\in \mathcal{S}} \omega_{\tilde{s}} ( \omega^* \cdot \hat{\Lambda}(X^*,\tilde{s}) )\right]
\end{equation}

Applying the Sherman-Morrison formula, we can invert the system above
\begin{equation}
    \hat{\Lambda}(X,s) =\left[ \gamma I - (\gamma - \psi^{-1}) \mathbb{1}_N \omega^*\right]  \tilde{b}^{\Lambda}(X,s).
\end{equation}

We can then write the expression above as follows:\small
\begin{equation}
    \frac{\hat{\Lambda}(X,s,s')}{\Lambda^*(X,s,s')} = 
    \gamma b^{\Lambda}(X,s,s') + (\gamma \psi - 1) \beta^* \left[ \omega^* \cdot \hat{\Lambda}(X^*,s') - \sum_{\tilde{s}\in \mathcal{S}} \omega_{\tilde{s}} ( \omega^* \cdot \hat{\Lambda}(X^*,\tilde{s}) )
    \right] -(\gamma - \psi^{-1}) \omega^* b^{\Lambda}(X,s).
\end{equation}\normalsize

\paragraph{Step 3: solving for the average $\hat{\Lambda}(X,s,s')$.} Assuming $X = X^*$, multiplying by $\omega_{s'}^*$, and adding across states, we obtain
\begin{equation}
    \omega^* \hat{\Lambda}(X^*,s) = \psi^{-1} \omega^* b^\Lambda(X^*,s).
\end{equation}

Averaging across $s$, we obtain
\begin{equation}
    \sum_{\tilde{s} \in \mathcal{S}} \omega^*_{\tilde{s}} [ \omega^* \hat{\Lambda}(X^*,\tilde{s}) ]= \psi^{-1} \sum_{\tilde{s} \in \mathcal{S}} \omega^*_{\tilde{s}} [\omega^* b^\Lambda(X^*,\tilde{s}) ].
\end{equation}

We can then write $\hat{\Lambda}(X,s,s')$ as follows\small
\begin{equation}
    \frac{\hat{\Lambda}(X,s,s')}{\Lambda^*(X,s,s')} = \gamma b^{\Lambda}(X,s,s') -(\gamma - \psi^{-1}) \omega^* b^{\Lambda}(X,s) + (\gamma  - \psi^{-1}) \beta^* \left[ \omega^* \cdot b^\Lambda(X^*,s') - \sum_{\tilde{s}\in \mathcal{S}} \omega_{\tilde{s}} ( \omega^* \cdot b^{\Lambda}(X^*,\tilde{s}) )
    \right].
\end{equation}\normalsize

\paragraph{Step 4: simplifying the expression for $b^\Lambda(X,s,s')$.} We can write $b^\Lambda(X,s,s')$ as follows:
\begin{align}
    b^\Lambda&(X,s,s')  = \frac{1}{\gamma}\left[\frac{\delta_{ss'}(X)}{p_{s'}^*} - \omega^* \cdot \delta_s(X)\right] + \frac{\beta^*}{\gamma} \left[ \omega^* \delta_{s'}(X) - \sum_{\tilde{s}} \omega^*_{\tilde{s}} \omega^* \cdot \delta_{\tilde{s}}(X) \right] \nonumber\\
&+ \frac{\psi-1}{1-\gamma} \left[ \omega^* \cdot \delta_{s}(X) - \beta^* \omega^* \cdot \delta_{s'}(X) + \beta^* \sum_{\tilde{s}} \omega_{\tilde{s}}^* (\omega^* \cdot \delta_{\tilde{s}}(X) ) \right]
- \left[ \log\frac{E^*}{E} - \left( \frac{E^*}{x_{s'}} - \frac{E}{x_s} \right) \right]\frac{\hat{\alpha}}{1+\nu}.
\end{align}

Combining terms, we obtain\footnotesize
\begin{align}
    &b^\Lambda(X,s,s')  = \frac{1}{\gamma}\frac{\delta_{ss'}(X)}{p_{s'}^*}- \frac{\psi - \gamma^{-1}}{\gamma-1} \left[\omega^* \cdot \delta_s(X) - \beta^* \omega^* \cdot \delta_{s'}(X)+\beta^* \sum_{\tilde{s}} \omega_{\tilde{s}}^* (\omega^* \cdot \delta_{\tilde{s}}(X) ) \right]
- \left[ \log\frac{E^*}{E} - \left( \frac{E^*}{x_{s'}} - \frac{E}{x_s} \right) \right]\frac{\hat{\alpha}}{1+\nu}.
\end{align}\normalsize

Notice that we can $\hat{\Lambda}(X,s,s')$ as follows:
\begin{equation}
    \frac{\hat{\Lambda}(X,s,s')}{\Lambda^*(X,s,s')} = \frac{\hat{\Lambda}(X^*,s,s')}{\Lambda^*(X,s,s')} + \psi^{-1} \left[ \log \frac{E}{E^*} - \frac{E- E^*}{x_s}\right] \frac{\hat{\alpha}}{1+\nu}.
\end{equation}

\end{proof}

\subsection{The economy with no labor frictions and iid returns}

Suppose that labor can be chosen conditional on the current productivity level. In this case, the problem of the firm can be written as
\begin{equation}
    \max_{h_{t}} x_t h_t^\alpha - w_t h_t,
\end{equation}
where $w_t \equiv \frac{W_t}{A_{t-1}}$. Labor demand takes the familiar form:
\begin{equation}
    \alpha x_t h_t^{\alpha-1} = w_t \Rightarrow h_t = \left( \frac{\alpha x_t}{w_t}\right)^{\frac{1}{1-\alpha}}.
\end{equation}

The labor supply from households is given by
\begin{equation}
    h_t = \left( \frac{w_t}{\xi} \right)^{\frac{1}{\nu}}.
\end{equation}

Combining labor supply and labor demand, we obtain the equilibrium hours and wages:
\begin{equation}
    h_t = \left( \frac{\alpha x_t}{\xi}\right)^{\frac{1}{1+\nu-\alpha}}, \qquad \qquad w_t = \xi^{\frac{1-\alpha}{1+\nu-\alpha}} (\alpha x_t)^{\frac{\nu}{1+ \nu-\alpha}}.
\end{equation}

Firm's profits are given by 
\begin{equation}
    \pi_t = A_{t-1} (1-\alpha) \left( \frac{\alpha}{\xi} \right)^{\frac{\alpha}{1+\nu-\alpha}} x_t^{\frac{1+\nu}{1+\nu-\alpha}}.
\end{equation}

Total surplus is given by
\begin{equation}
    \tilde{C}_t = A_{t-1} \left[ x_t h_t^\alpha - \xi \frac{h_{t}^{1+\nu}}{1+\nu} \right] = A_{t-1} \left(1 - \frac{\alpha}{1+\nu} \right) \left( \frac{\alpha}{\xi} \right)^{\frac{\alpha}{1+\nu-\alpha}} x_t^{\frac{1+\nu}{1+\nu-\alpha}}.
\end{equation}

Let $P(X,s)$ denote the price of the surplus claim normalized by lagged productivity. Market clearing condition implies that
\begin{equation}
    1-\beta^* = \left(1 - \frac{\alpha}{1+\nu} \right) \left( \frac{\alpha}{\xi} \right)^{\frac{\alpha}{1+\nu-\alpha}} \frac{x_s^{\frac{1+\nu}{1+\nu-\alpha}}}{P(X,s)},
\end{equation}
where $1-\beta^*$ is the consumption-wealth ratio, which we assume to be constant.

The return on the surplus claim is given by
\begin{equation}
    R_p(X,s,s') = \frac{x_s}{\beta^*} \frac{x_{s'}^{\frac{1+\nu}{1+\nu-\alpha}}}{x_s^{\frac{1+\nu}{1+\nu-\alpha}}} = \frac{x_{s'}}{\beta^*} \left( \frac{x_{s'}}{x_s} \right)^{ \frac{\alpha}{1+\nu-\alpha} }.
\end{equation}











% \subsection{Perturbation}

% Consider a family of economies indexed by $\epsilon$, where the EIS and coefficient of relative risk aversion are given by $\psi = 1 + \hat{\psi} \epsilon$ and $\gamma = 1 + \hat{\gamma} \epsilon$. 
% Therefore, we obtain the log case when $\epsilon = 0$ and we allow for small deviations away from unit EIS and unit risk aversion when $\epsilon > 0$. 

% All equilibrium objects are now indexed by $\epsilon$, for instance, the net worth multiplier is now given by $v(X,s; \epsilon)$. We are interested in the expansion of $v(X,s; \epsilon)$ on $\epsilon$, for $\epsilon$ small:
% \begin{equation}
%     v(X,s;\epsilon) = v^*(X,s) + \tilde{v}(X,s) \epsilon + \mathcal{O}(\epsilon^2),
% \end{equation}
% where $v^*(X,s) \equiv v(X,s; 0)$ and $\tilde{v}(X,s)$ represents the first-order correction of $v(X,s;\epsilon)$ in $\epsilon$. 

% Similarly, we can write the consumption-wealth ratio  $c(X,s;\epsilon)$ as follows:
% \begin{equation}
%     c(X,s;\epsilon) = c^*(X,s) + \tilde{c}(X,s) \epsilon + \mathcal{O}(\epsilon^2),
% \end{equation}
% and analogously for the remaining equilibrium variables. 

% The functions $v^*(X,s)$ and $c^*(X,s)$ are already known, as they correspond to the solution in the log case, which we characterized above. It remains to solve for $\tilde{v}(X,s)$,  $\tilde{c}(X,s)$, and the first-order correction for the other variables.

% From the Bellman equation for household $i$, we can write $\hat{v}_i(X,s) \equiv \log v_i(X,s)$ as follows:
% \begin{equation}
%     \hat{v}_i = \frac{1}{1-\psi^{-1}} \log \left[(1-\beta) \exp[(1-\psi^{-1}) \hat{c}_i] + \beta \exp[ (1-\psi^{-1}) \mathcal{Q}(\hat{v}_i' + \hat{n}_i')]
%      \right],
% \end{equation}
% where
% \begin{equation}
%     \mathcal{Q}(z') = \frac{1}{1-\gamma} \log \mathbb{E}_i\left[ \exp [(1-\gamma) z']\right]
% \end{equation}

% For $\epsilon = 0$, we obtain
% \begin{equation}
%     \hat{v}_i^*(X,s) = (1-\beta) \hat{c}_i^*(X,s) + \beta \mathbb{E}_i[\hat{v}_i^*(X',s') + \hat{n}_i^*(X,s,s')].
% \end{equation}

% We follow \citet{hansen2007intertemporal} and expand the logarithm and exponential functions in the recursion for $\hat{v}_i$ and we will keep up to second-order terms in $\hat{c}_i$ and $\mathcal{Q}(\cdot)$. We can use the approximate recursion to compute the derivative of $\hat{v}_i$ with respect to $\epsilon$ at $\epsilon = 0$. Considering higher order terms or even the exact recursion for $\hat{v}_i$ would yield the same result for the derivative at the cost of more cumbersome algebra. The expansion of the recursion is given by
% \begin{align*}
%     \hat{v}_i &= \frac{1}{1-\psi^{-1}} \log \left[ 1 + (1-\beta) (\exp[(1-\psi^{-1}) \hat{c}_i]-1) + \beta (\exp[ (1-\psi^{-1}) \mathcal{Q}(\hat{v}_i' + \hat{n}_i')]-1)
%      \right],\\
%      &\approx\frac{1}{1-\psi^{-1}} \log \left[ 1 + (1-\beta) (1-\psi^{-1})\left( \hat{c}_i + (1-\psi^{-1}) \frac{\hat{c}_i^2}{2} \right) + \beta (1-\psi^{-1}) \left( \mathcal{Q} + (1-\psi^{-1}) \frac{\mathcal{Q}^2}{2} \right)
%      \right],\\
%      &\approx  (1-\beta) \left( \hat{c}_i + (1-\psi^{-1}) \frac{\hat{c}_i^2}{2} \right) + \beta \left( \mathcal{Q} + (1-\psi^{-1}) \frac{\mathcal{Q}^2}{2} \right) - \frac{1-\psi^{-1}}{2} \left( (1-\beta) \hat{c}_i + \beta \mathcal{Q}\right)^2\\
%      &= (1-\beta) \hat{c}_i + \beta \mathcal{Q} + \beta (1-\beta) (1-\psi^{-1}) \frac{\hat{c}_i^2-2 \hat{c}_i \mathcal{Q} + \mathcal{Q}^2}{2}.
% \end{align*}

% The approximate recursion is then given by
% \begin{equation}
%     \hat{v}_i \approx (1-\beta) \hat{c}_i + \beta \mathcal{Q} + \beta (1-\beta) (1-\psi^{-1}) \frac{(\hat{c}_i- \mathcal{Q})^2}{2}.
% \end{equation}

% Similarly, we can write $\mathcal{Q}$ as follows
% \begin{align}
%     \mathcal{Q}(\hat{v}_i' + \hat{n}_i') &= \frac{1}{1-\gamma} \log \mathbb{E}_i\left[1 + \exp[(1-\gamma) (\hat{v}_i' + \hat{n}_i') -1\right]\\
%     &\approx \frac{1}{1-\gamma} \log \mathbb{E}_i\left[1 + (1-\gamma) (\hat{v}_i' + \hat{n}_i') + (1-\gamma)^2 \frac{(\hat{v}_i' + \hat{n}_i')^2}{2}\right]\\
%     &\approx \mathbb{E}_i\left[\hat{v}_i' + \hat{n}_i' + (1-\gamma) \frac{(\hat{v}_i' + \hat{n}_i')^2}{2} - \frac{(1-\gamma)}{2} (\hat{v}_i' + \hat{n}_i')^2\right].
% \end{align}

